\documentclass[11pt,a4paper]{article}

\usepackage[T1]{fontenc}
\usepackage[utf8]{inputenc}
\usepackage{lmodern}
\usepackage[margin=1in]{geometry}
\usepackage{amsmath,amssymb}
\usepackage{enumitem}
\usepackage{booktabs}
\usepackage{tabularx}
\usepackage{longtable}
\usepackage{verbatim}
\usepackage{fancyvrb}
\usepackage[dvipsnames]{xcolor}
\usepackage{framed}
\usepackage{newunicodechar}
\usepackage{hyperref}
\hypersetup{
  colorlinks=true,
  urlcolor=blue!50!black,
  linkcolor=blue!50!black,
  citecolor=blue!50!black,
  pdftitle={Sentrix: A Layer-1 Blockchain --- Protocol Specification},
  pdfauthor={Satya Kwok}
}
\usepackage{titlesec}
\titlespacing{\section}{0pt}{14pt}{6pt}
\titlespacing{\subsection}{0pt}{10pt}{4pt}

\setlength{\parskip}{6pt}
\setlength{\parindent}{0pt}

% Map the unicode glyphs we use in inline code to LaTeX math/text equivalents
% so pdflatex (no XeTeX/LuaTeX) can typeset them.
\newunicodechar{⌊}{\ensuremath{\lfloor}}
\newunicodechar{⌋}{\ensuremath{\rfloor}}
\newunicodechar{⌈}{\ensuremath{\lceil}}
\newunicodechar{⌉}{\ensuremath{\rceil}}
\newunicodechar{−}{-}
\newunicodechar{≥}{\ensuremath{\geq}}
\newunicodechar{≤}{\ensuremath{\leq}}
\newunicodechar{≠}{\ensuremath{\neq}}
\newunicodechar{×}{\ensuremath{\times}}
\newunicodechar{÷}{\ensuremath{\div}}
\newunicodechar{∈}{\ensuremath{\in}}
\newunicodechar{∀}{\ensuremath{\forall}}
\newunicodechar{∃}{\ensuremath{\exists}}
\newunicodechar{Σ}{\ensuremath{\Sigma}}
\newunicodechar{Δ}{\ensuremath{\Delta}}
\newunicodechar{∎}{\ensuremath{\blacksquare}}
\newunicodechar{→}{\ensuremath{\rightarrow}}
\newunicodechar{←}{\ensuremath{\leftarrow}}
\newunicodechar{↔}{\ensuremath{\leftrightarrow}}
\newunicodechar{⇒}{\ensuremath{\Rightarrow}}
\newunicodechar{⊥}{\ensuremath{\bot}}
\newunicodechar{′}{\ensuremath{\prime}}
\newunicodechar{″}{\ensuremath{\prime\prime}}
\newunicodechar{∪}{\ensuremath{\cup}}
\newunicodechar{∩}{\ensuremath{\cap}}
\newunicodechar{·}{\ensuremath{\cdot}}
\newunicodechar{²}{\ensuremath{^2}}
\newunicodechar{³}{\ensuremath{^3}}
\newunicodechar{⁰}{\ensuremath{^0}}
\newunicodechar{¹}{\ensuremath{^1}}
\newunicodechar{⁴}{\ensuremath{^4}}
\newunicodechar{⁵}{\ensuremath{^5}}
\newunicodechar{⁶}{\ensuremath{^6}}
\newunicodechar{⁸}{\ensuremath{^8}}
\newunicodechar{⁻}{\ensuremath{^-}}
\newunicodechar{ₐ}{\ensuremath{_a}}
\newunicodechar{₀}{\ensuremath{_0}}
\newunicodechar{₁}{\ensuremath{_1}}
\newunicodechar{₂}{\ensuremath{_2}}
\newunicodechar{ₕ}{\ensuremath{_h}}
\newunicodechar{β}{\ensuremath{\beta}}
\newunicodechar{α}{\ensuremath{\alpha}}
\newunicodechar{ε}{\ensuremath{\varepsilon}}
\newunicodechar{𝒮}{\ensuremath{\mathcal{S}}}
\newunicodechar{ℬ}{\ensuremath{\mathcal{B}}}
\newunicodechar{|}{|}
\newunicodechar{ }{ }

% Pandoc helpers
\providecommand{\passthrough}[1]{#1}
\providecommand{\tightlist}{\setlength{\itemsep}{0pt}\setlength{\parskip}{0pt}}

% Pandoc's syntax-highlighting environment (used for code blocks)
\providecommand{\NormalTok}[1]{#1}
\providecommand{\KeywordTok}[1]{\textbf{#1}}
\providecommand{\DataTypeTok}[1]{\textbf{#1}}
\providecommand{\CommentTok}[1]{\textit{#1}}
\providecommand{\StringTok}[1]{#1}

\definecolor{shadecolor}{rgb}{0.95,0.95,0.95}

% Highlighting/Shaded environments pandoc emits
\newenvironment{Shaded}{}{}
\newenvironment{Highlighting}[1][]{\par\noindent\begin{minipage}{\linewidth}\ttfamily\small}{\end{minipage}\par}

\title{\textbf{Sentrix}\\[2pt]\large A Layer-1 Blockchain --- Protocol Specification}
\author{Satya Kwok \\ \texttt{satya@sentrixchain.com} \\ \texttt{sentrixchain.com}}
\date{Version 1.3.0}

\begin{document}
\maketitle
\tableofcontents
\newpage

\hypertarget{sentrix}{%
\section{Sentrix}\label{sentrix}}

\hypertarget{a-layer-1-blockchain--protocol-specification}{%
\subsubsection{A Layer-1 Blockchain --- Protocol
Specification}\label{a-layer-1-blockchain--protocol-specification}}

\textbf{Author:} Satya Kwok
\textless{}\href{mailto:satya@sentrixchain.com}{\nolinkurl{satya@sentrixchain.com}}\textgreater{}
\textbf{Web:} sentrixchain.com \textbf{Version:} 1.3.0 (supersedes
1.2.4)

\begin{center}\rule{0.5\linewidth}{0.5pt}\end{center}

\hypertarget{abstract}{%
\subsection{Abstract}\label{abstract}}

This document specifies \emph{Sentrix}, a Layer-1 blockchain that
combines a Tendermint-style Byzantine Fault Tolerant (BFT) consensus
protocol with delegated proof-of-stake validator selection, a
deterministic dual-rail execution layer (a native protocol-operation
rail and an Ethereum Virtual Machine rail running on \texttt{revm}), and
a binary sparse Merkle tree state commitment. We give a formal systems
model, the consensus protocol\textquotesingle s message structure and
round logic with safety/liveness conditions, the execution pipeline
including a forward-looking parallel-execution design, a performance
model with closed-form throughput and latency expressions, an
adversarial model with quantitative security bounds, and explicit
failure-handling protocols for the four classes of fault that the
production deployment has encountered. The paper is intended as an
engineer-auditable specification rather than a marketing document.

The chain runs in production on chain ID \texttt{7119} (mainnet) and
\texttt{7120} (testnet) under the open-source Rust implementation at
\texttt{github.com/sentrix-labs/sentrix}. Economic constants --- 315M
SRX maximum supply, 1 SRX initial block reward halving every 126M
blocks, 50\% of every transaction fee burned --- are codified at the
protocol level and are not subject to governance.

\begin{center}\rule{0.5\linewidth}{0.5pt}\end{center}

\hypertarget{1-introduction}{%
\subsection{1. Introduction}\label{1-introduction}}

\hypertarget{11-scope}{%
\subsubsection{1.1 Scope}\label{11-scope}}

This document specifies the Sentrix protocol at a level sufficient for
an engineer to (a) implement a conforming node, (b) audit a divergence
between two implementations, or (c) reason about chain behaviour under
adversarial or partial-failure conditions. It is not an introduction to
distributed systems or to blockchain primitives; familiarity with
state-machine replication, Byzantine consensus, and the Ethereum
execution model is assumed.

Constants, fork heights, slashing parameters, and routing addresses
cited throughout this document are pinned to the Rust implementation in
the workspace at \texttt{github.com/sentrix-labs/sentrix} as of release
\texttt{v2.1.56}. Where the implementation has parametric values
configured at genesis or by environment variable, the document marks the
value as \emph{parametric} and points at the relevant module.

\hypertarget{12-design-goals}{%
\subsubsection{1.2 Design Goals}\label{12-design-goals}}

The protocol prioritizes, in order:

\begin{enumerate}
\def\labelenumi{\arabic{enumi}.}
\tightlist
\item
  \textbf{Deterministic state agreement} --- every honest node executing
  the same blocks against the same prior state produces a bit-identical
  result.
\item
  \textbf{Single-block finality} --- once a block is committed, no fork
  can replace it without violating the BFT safety threshold.
\item
  \textbf{EVM compatibility} --- Solidity contracts and standard
  Ethereum tooling work without modification.
\item
  \textbf{One-second target block time} --- at the BFT round structure
  described in §4 and under the network assumptions of §2.2.
\item
  \textbf{Operator simplicity} --- the chain ships as a single Rust
  binary; an operator runs one process and manages one local key-value
  store.
\item
  \textbf{Predictable monetary policy} --- supply, halving cadence, fee
  burn ratio, and genesis allocation are not governable.
\end{enumerate}

\hypertarget{13-non-goals}{%
\subsubsection{1.3 Non-Goals}\label{13-non-goals}}

The protocol explicitly does not target:

\begin{itemize}
\tightlist
\item
  \textbf{Maximum throughput at any cost.} The throughput bound (§10)
  trades off against single-block finality and small validator-set
  requirements.
\item
  \textbf{On-chain privacy.} Every transaction, balance, and call is
  publicly observable. Privacy is delegated to the application layer.
\item
  \textbf{Sharding.} State and execution are unsharded. Section 10
  derives the consequent scalability bounds.
\item
  \textbf{Multi-VM polyglot execution.} The protocol natively supports
  the EVM and a fixed-ABI native rail; arbitrary VM plug-ins are out of
  scope.
\end{itemize}

\begin{center}\rule{0.5\linewidth}{0.5pt}\end{center}

\hypertarget{2-system-model}{%
\subsection{2. System Model}\label{2-system-model}}

\hypertarget{21-notation}{%
\subsubsection{2.1 Notation}\label{21-notation}}

Throughout, we use the following symbols:

\begin{longtable}[]{@{}ll@{}}
\toprule\noalign{}
Symbol & Meaning \\
\midrule\noalign{}
\endhead
\bottomrule\noalign{}
\endlastfoot
\texttt{n} & Number of active validators in the current epoch \\
\texttt{f} & Maximum number of Byzantine validators tolerated,
\texttt{f\ =\ ⌊(n\ −\ 1)\ /\ 3⌋} \\
\texttt{Q} & BFT supermajority quorum, \texttt{Q\ =\ ⌊2n/3⌋\ +\ 1}. For
\texttt{n\ =\ 4}, \texttt{Q\ =\ 3}. For \texttt{n\ =\ 21},
\texttt{Q\ =\ 15}. \\
\texttt{S} & Canonical chain state, \texttt{S\ ∈\ 𝒮} \\
\texttt{S₀} & Genesis state \\
\texttt{B} & Block, \texttt{B\ =\ (header,\ txs)} \\
\texttt{H(B)} & Cryptographic hash of \texttt{B}\textquotesingle s
canonical encoding \\
\texttt{STF} & State transition function,
\texttt{STF:\ 𝒮\ ×\ ℬ\ →\ 𝒮\ ∪\ \{⊥\}} \\
\texttt{h} & Block height, monotonically increasing \\
\texttt{r} & BFT round at a given height,
\texttt{r\ ∈\ \{0,\ 1,\ …,\ MAX\_ROUND\}} \\
\texttt{Δ} & Upper bound on message delay after global stabilisation
time (GST) \\
\texttt{Vₐ} & Active validator set for the current epoch, ` \\
\texttt{s(v)} & Stake-weight of validator \texttt{v\ ∈\ Vₐ} \\
\end{longtable}

\texttt{SRX} is the native chain asset; \texttt{1\ SRX\ =\ 10⁸\ sentri}
(the smallest indivisible unit). All token amounts in this document are
expressed in sentri unless explicitly noted otherwise.

\hypertarget{22-network-model}{%
\subsubsection{2.2 Network Model}\label{22-network-model}}

We assume \textbf{partial synchrony} in the sense of
Dwork--Lynch--Stockmeyer {[}15{]}: there exists an unknown but finite
time \texttt{GST} (Global Stabilisation Time) after which all messages
between honest participants are delivered within bounded delay
\texttt{Δ}. Before \texttt{GST}, messages may be arbitrarily delayed,
reordered, or dropped. The protocol must remain \emph{safe} under
arbitrary asynchrony and become \emph{live} once partial synchrony
holds.

Inter-validator communication is over a libp2p {[}11{]} mesh (§7) with
three message classes (block gossip, transaction gossip, BFT messages).
Messages are authenticated by validator signatures. The protocol does
not assume FIFO delivery or reliable transport; the application layer is
responsible for retransmission and deduplication.

\hypertarget{23-adversarial-model}{%
\subsubsection{2.3 Adversarial Model}\label{23-adversarial-model}}

We adopt the standard BFT adversarial model:

\begin{itemize}
\tightlist
\item
  \textbf{Honest validators} follow the protocol exactly.
\item
  \textbf{Byzantine validators} may deviate arbitrarily --- including
  signing conflicting messages, withholding messages, equivocating on
  proposals, or coordinating with other Byzantine validators.
\item
  The adversary controls a stake-weighted fraction \texttt{β} of the
  active set, \texttt{β\ ∈\ {[}0,\ 1{]}}.
\end{itemize}

The protocol guarantees safety provided \texttt{β\ \textless{}\ 1/3}
(i.e. fewer than \texttt{n/3} validators by stake are Byzantine). It
guarantees liveness under the additional partial-synchrony assumption
(§2.2) when \texttt{β\ \textless{}\ 1/3}.

The adversary is computationally bounded; specifically, the adversary
cannot forge ECDSA signatures, find SHA-256 collisions, or break the
secp256k1 group\textquotesingle s discrete-log assumption.

\hypertarget{24-state-machine-replication}{%
\subsubsection{2.4 State Machine
Replication}\label{24-state-machine-replication}}

Sentrix is an instance of state machine replication. Each validator
maintains an identical replica of \texttt{S}. Clients submit
transactions; consensus orders them into blocks; every replica applies
blocks via \texttt{STF} in the same order, deriving the same successor
states.

Formally:

\begin{quote}
\textbf{(SMR Property)} For every honest validator \texttt{v} and every
height \texttt{h}, after \texttt{v} has applied all blocks
\texttt{B₁,\ …,\ Bₕ}, the resulting state \texttt{Sₕ} is independent of
\texttt{v}. That is, \texttt{∀\ v,\ v′\ honest:\ Sₕ(v)\ =\ Sₕ(v′)}.
\end{quote}

The SMR property is achieved by (a) the consensus protocol agreeing on a
single block per height (§4 safety) and (b) the determinism of
\texttt{STF} (§5.2).

\begin{center}\rule{0.5\linewidth}{0.5pt}\end{center}

\hypertarget{3-architecture}{%
\subsection{3. Architecture}\label{3-architecture}}

The node is a single Rust binary composed of subsystems sharing the
canonical state \texttt{S}:

\begin{Shaded}
\begin{verbatim}
  Client / Wallet --tx--> sentrix-rpc --admit--> mempool
                                                    |
                                       drain on propose v
                                                  sentrix-bft
                                                <- gossip ->   sentrix-network <- peers ->
                                                    |
                                          FinalizeBlock v
                                              block_executor
                                              |           |
                                       Native rail    sentrix-evm (revm 37)
                                              |           |
                                              v           v
                                              sentrix-trie (binary SMT, 256-level)
                                                          |
                                                          v
                                              sentrix-storage (MDBX KV) --> chain.db
\end{verbatim}
\end{Shaded}

\textbf{Figure 1.} System architecture. The block executor receives a
\texttt{FinalizeBlock} action from the BFT engine after consensus and
dispatches each transaction via the routing rules in §5.5. The trie +
storage layer is shared across rails; there is no separate execution
database.

The 14 crates of the workspace
(\texttt{crates/sentrix-\{primitives,codec,wire,trie,storage,wallet,bft,staking,network,evm,precompiles,core,rpc-types,rpc\}})
compile into the single \texttt{sentrix} binary in \texttt{bin/sentrix}.
A separate \texttt{bin/sentrix-faucet} exists for testnet operations and
is not part of consensus.

\begin{center}\rule{0.5\linewidth}{0.5pt}\end{center}

\hypertarget{4-consensus-protocol-voyager-bft}{%
\subsection{4. Consensus Protocol (Voyager
BFT)}\label{4-consensus-protocol-voyager-bft}}

Sentrix consensus is a Tendermint-derivative we call \emph{Voyager}.
This section specifies the round structure, message types, locking
rules, and safety/liveness conditions.

\hypertarget{41-validator-selection-dpos}{%
\subsubsection{4.1 Validator Selection
(DPoS)}\label{41-validator-selection-dpos}}

The active validator set \texttt{Vₐ} for an epoch is selected by
stake-weighted ranking. Any account that has executed
\texttt{RegisterValidator} and bonded at least the genesis-configured
minimum self-stake is a \emph{candidate}. Token holders may delegate
stake to candidates via the \texttt{Delegate} operation. The active set
for epoch \texttt{e\ +\ 1} is the
top-\texttt{MAX\_ACTIVE\_VALIDATORS\ =\ 21} candidates ranked by
\texttt{total\_stake\ =\ self\_stake\ +\ total\_delegated}, computed
deterministically at the boundary block
\texttt{h\ =\ e\ ×\ EPOCH\_LENGTH} where
\texttt{EPOCH\_LENGTH\ =\ 28\_800} blocks (\textasciitilde8 hours at 1 s
blocks).

Stake withdrawal initiates an \emph{unbonding period} of
\texttt{UNBONDING\_PERIOD\ =\ 201\_600} blocks, during which the stake
remains slashable but does not earn rewards. This prevents a Byzantine
validator from offloading slashable stake immediately before
equivocating.

\texttt{AddSelfStake} (gated by \texttt{ADD\_SELF\_STAKE\_HEIGHT},
activated 2026-04-28) lets a registered validator increase its
self-stake from its own wallet without re-registering, used during
unjail flows and stake top-ups.

\hypertarget{42-round-structure}{%
\subsubsection{4.2 Round Structure}\label{42-round-structure}}

A \emph{height} \texttt{h} is finalized through one or more
\emph{rounds} \texttt{r\ ∈\ \{0,\ 1,\ …,\ MAX\_ROUND\}} where
\texttt{MAX\_ROUND\ =\ 100}. Each round consists of three phases ---
\texttt{PROPOSE}, \texttt{PREVOTE}, \texttt{PRECOMMIT} --- followed by a
\texttt{FINALIZE} phase if quorum is reached. The four phases are
encoded in
\texttt{BftPhase\ =\ \{Propose,\ Prevote,\ Precommit,\ Finalize\}}.

For round \texttt{(h,\ r)}, the \emph{proposer} is chosen by

\[
\text{propose}(h, r) = V_a[(h + r) \bmod n]
\]

Round-robin over the active set; stake-weight does not affect proposer
selection within an epoch (the function name \texttt{weighted\_proposer}
in \texttt{crates/sentrix-staking/src/staking.rs} is historical).

Three message types circulate:

\begin{longtable}[]{@{}lll@{}}
\toprule\noalign{}
Type & Carried fields & Role \\
\midrule\noalign{}
\endhead
\bottomrule\noalign{}
\endlastfoot
\texttt{Proposal(h,\ r,\ B)} &
\texttt{height,\ round,\ block\_hash,\ sig} & Proposer\textquotesingle s
claim that \texttt{B} is the next block at \texttt{(h,\ r)} \\
\texttt{Prevote(h,\ r,\ x)} &
\texttt{height,\ round,\ x\ ∈\ \{block\_hash,\ ⊥\},\ sig} &
Validator\textquotesingle s vote \\
\texttt{Precommit(h,\ r,\ x)} &
\texttt{height,\ round,\ x\ ∈\ \{block\_hash,\ ⊥\},\ sig} &
Validator\textquotesingle s commitment \\
\end{longtable}

A signature is a 65-byte secp256k1 ECDSA recoverable signature (64-byte
compact \texttt{(r,\ s)} plus one byte \texttt{recovery\_id}), computed
over the SHA-256 of a canonical signing payload that includes a
domain-separation prefix (\texttt{0x01} proposal, \texttt{0x02} prevote,
\texttt{0x03} precommit). The proposer\textquotesingle s signature on a
\texttt{Proposal} also covers \texttt{B}\textquotesingle s contents.
Wire-level messages are wrapped in
\texttt{BftMessage\ =\ \{Propose(Proposal),\ Prevote(Prevote),\ Precommit(Precommit),\ RoundStatus(RoundStatus)\}}
for transport.

\hypertarget{421-propose}{%
\paragraph{4.2.1 PROPOSE}\label{421-propose}}

The proposer for \texttt{(h,\ r)} constructs a candidate block
\texttt{B}:

\begin{verbatim}
B = {
    header: {
        index: h,
        previous_hash: H(B_{h-1}),
        timestamp: now(),
        proposer: addr(propose(h, r)),
        state_root: STF(S_{h-1}, B).root,    // STAGED — see §5
        justification: J_{h-1},               // precommits proving B_{h-1} finalized
        merkle_root: tx_merkle(B.txs),
    },
    txs: drain_mempool(MAX_TX_PER_BLOCK = 5_000)
}
\end{verbatim}

The proposer broadcasts \texttt{Proposal(h,\ r,\ B)} to all of
\texttt{Vₐ}. Honest validators wait at most
\texttt{propose\_timeout(r)\ =\ min(20\_000\ +\ r\ ×\ 1\_000,\ 30\_000)\ ms}
for the proposal; on timeout they treat the proposal as \texttt{⊥} and
proceed to \texttt{PREVOTE}.

These timeouts are upper bounds, not phase durations. In the happy path
on a low-latency mesh, all three phases complete in well under one
second; the chain produces blocks at the
\texttt{BLOCK\_TIME\_SECS\ =\ 1} target. The wide timeouts exist to
absorb peer-mesh repair after transient disconnects (see
\texttt{crates/sentrix-bft/src/engine/timeouts.rs} for the incident
history).

\hypertarget{422-prevote}{%
\paragraph{4.2.2 PREVOTE}\label{422-prevote}}

Each validator \texttt{v\ ∈\ Vₐ} evaluates the proposal:

\begin{itemize}
\tightlist
\item
  If \texttt{B} validates (signature, parent hash, transaction
  signatures, state root) and \texttt{v} is not locked on a different
  block at this height, \texttt{v} broadcasts
  \texttt{Prevote(h,\ r,\ H(B))}.
\item
  Otherwise, \texttt{v} broadcasts \texttt{Prevote(h,\ r,\ ⊥)}.
\end{itemize}

The validator then waits for prevotes from validators carrying total
stake ≥ \texttt{Q}. If \texttt{Q} validators by stake-weight prevote the
same \texttt{H(B)} (a \emph{polka}), \texttt{v} proceeds to
\texttt{PRECOMMIT} with \texttt{H(B)}. Otherwise \texttt{v} proceeds
with \texttt{⊥}.

Timeout:
\texttt{prevote\_timeout(r)\ =\ min(12\_000\ +\ r\ ×\ 2\_000,\ 30\_000)\ ms}.

\hypertarget{423-precommit}{%
\paragraph{4.2.3 PRECOMMIT}\label{423-precommit}}

Each validator broadcasts \texttt{Precommit(h,\ r,\ x)}:

\begin{itemize}
\tightlist
\item
  \texttt{x\ =\ H(B)} if \texttt{v} saw a polka for \texttt{H(B)} in the
  prevote phase.
\item
  \texttt{x\ =\ ⊥} otherwise.
\end{itemize}

A validator waits for precommits carrying total stake ≥ \texttt{Q}. If
\texttt{Q} precommits agree on \texttt{H(B)}, \texttt{B} is
\emph{finalized} at height \texttt{h}.

Timeout:
\texttt{precommit\_timeout(r)\ =\ min(12\_000\ +\ r\ ×\ 2\_000,\ 30\_000)\ ms}.

\hypertarget{424-finalize}{%
\paragraph{4.2.4 FINALIZE}\label{424-finalize}}

The set of \texttt{Q} precommits constitutes the \emph{justification}
\texttt{Jₕ} for block \texttt{B}. \texttt{Jₕ} is included in the next
block\textquotesingle s header
(\texttt{B\_\{h+1\}.header.justification}), providing public proof of
finalization.

Each validator applies \texttt{B} to its local state via \texttt{STF}
(§5) and advances to round \texttt{0} of height \texttt{h\ +\ 1}.

\hypertarget{43-locking}{%
\subsubsection{4.3 Locking}\label{43-locking}}

Once a validator has prevoted a block in round \texttt{r}, it locks on
that block and round (\texttt{locked\_block} and \texttt{locked\_round}
are persisted in \texttt{BftRoundState}). In subsequent rounds at the
same height, it will only prevote a different block if it observes a
polka (≥ \texttt{Q} prevotes for the same alternative) at a
\emph{higher} round. This rule is what guarantees safety:

\begin{quote}
\textbf{(Locking Invariant)} An honest validator that precommits
\texttt{H(B)} in round \texttt{r} only precommits a different block
\texttt{H(B′)} in round \texttt{r′\ \textgreater{}\ r} if \texttt{≥\ Q}
prevotes for \texttt{H(B′)} are observed at some round \texttt{r′′} with
\texttt{r\ ≤\ r′′\ \textless{}\ r′}.
\end{quote}

Since \texttt{≥\ Q} honest validators must have abandoned their lock on
\texttt{H(B)} to prevote \texttt{H(B′)}, and abandoning a lock requires
observing the alternative polka, no honest validator precommits both
\texttt{H(B)} and \texttt{H(B′)}.

\hypertarget{44-round-skip}{%
\subsubsection{4.4 Round Skip}\label{44-round-skip}}

If a round fails to finalize within its phase timeouts plus \texttt{ε},
the round terminates and round \texttt{r\ +\ 1} begins. The proposer
rotates to \texttt{Vₐ{[}(h\ +\ r\ +\ 1)\ mod\ n{]}}.

Round advancement is \textbf{timeout-only} (per \texttt{2026-04-17} fix,
see \texttt{crates/sentrix-bft/src/lib.rs} doc): no vote-triggered or
\texttt{RoundStatus}-triggered catch-up promotes a validator across
rounds. This avoids the \emph{validator-leapfrog} stall in which
validators clear collected votes on every round jump.

Per-phase timeouts grow linearly with round number to absorb sustained
partial synchrony:

\begin{verbatim}
propose_timeout(r)   = min(20_000 + r × 1_000, 30_000) ms
prevote_timeout(r)   = min(12_000 + r × 2_000, 30_000) ms
precommit_timeout(r) = min(12_000 + r × 2_000, 30_000) ms
\end{verbatim}

Once any phase timeout reaches the cap of 30 s, it does not grow
further. After \texttt{MAX\_ROUND\ =\ 100} consecutive failed rounds at
a single height, the chain is considered stalled at that height;
recovery requires operator intervention (§11.1).

\hypertarget{45-consensus-sequence}{%
\subsubsection{4.5 Consensus Sequence}\label{45-consensus-sequence}}

\begin{Shaded}
\begin{verbatim}
PROPOSE phase     PREVOTE phase     PRECOMMIT phase    FINALIZE
(timeout 20s+r*1s) (timeout 12s+r*2s) (timeout 12s+r*2s) (apply STF)
caps 30s          caps 30s          caps 30s
     |                 |                 |                 |
     v                 v                 v                 v
  proposer P     each validator    each validator    block B finalized
  for height H,  prevotes B if     precommits B if    J_h = {Precommits}
  round R        valid             >=Q stake-weight    advance to (h+1, 0)
  proposes B                       prevotes for B
\end{verbatim}
\end{Shaded}

\textbf{Figure 2.} Voyager BFT round in the happy path. The labelled
timeouts are upper bounds on each phase; on a healthy mesh the round
completes in much less time, and the chain produces blocks at the 1 s
target.

\hypertarget{46-safety}{%
\subsubsection{4.6 Safety}\label{46-safety}}

\textbf{Theorem 1 (Agreement).} \emph{For all heights \texttt{h}, no two
honest validators finalize different blocks at height \texttt{h},
provided \texttt{β\ \textless{}\ 1/3} of stake is Byzantine.}

\textbf{Proof sketch.} Suppose honest validators \texttt{v₁} and
\texttt{v₂} finalize \texttt{B} and \texttt{B′} at height \texttt{h}
with \texttt{H(B)\ ≠\ H(B′)}. Each saw \texttt{Q} precommits for their
respective blocks. With \texttt{Q\ =\ ⌊2n/3⌋\ +\ 1} and
\texttt{f\ =\ ⌊(n\ −\ 1)/3⌋}, we have \texttt{2Q\ −\ n\ ≥\ f\ +\ 1}, so
the two precommit sets share at least \texttt{f\ +\ 1} validators. Among
these shared validators at least one is honest (since at most \texttt{f}
are Byzantine). An honest validator that precommits \texttt{H(B)} cannot
precommit \texttt{H(B′)} without an intervening polka for \texttt{H(B′)}
(Locking Invariant, §4.3). For the polka to occur, validators carrying
total stake \texttt{≥\ Q} must have prevoted \texttt{H(B′)}, which
requires at least \texttt{Q\ −\ f\ \textgreater{}\ f} honest validators
to have abandoned their lock on \texttt{H(B)} --- contradicting the
Locking Invariant for those validators. ∎

\textbf{Theorem 2 (Validity).} \emph{Every finalized block is
well-formed: its header validates, all transaction signatures verify,
the state root matches \texttt{STF} execution, and its justification is
a valid \texttt{Q}-precommit set on its parent.}

\textbf{Proof.} Honest validators only prevote a well-formed block
(§4.2.2). For finalization, \texttt{Q\ ≥\ n\ −\ f} validators precommit,
of which at least \texttt{Q\ −\ f\ ≥\ n\ −\ 2f\ \textgreater{}\ f} are
honest. Honest precommits require honest prevotes, and honest prevotes
require valid blocks. ∎

\hypertarget{47-liveness}{%
\subsubsection{4.7 Liveness}\label{47-liveness}}

\textbf{Theorem 3 (Termination).} \emph{Under partial synchrony (§2.2)
and \texttt{β\ \textless{}\ 1/3}, every height eventually finalizes
within \texttt{MAX\_ROUND} rounds.}

\textbf{Proof sketch.} After \texttt{GST}, all messages between honest
validators deliver within \texttt{Δ}. The round-skip mechanism (§4.4)
eventually selects a timeout greater than \texttt{Δ} (capped at 30 s,
sufficient for any reasonable \texttt{Δ} in production network
conditions). Within such a round, an honest proposer\textquotesingle s
\texttt{Proposal} reaches all honest validators in time, all honest
prevotes reach all honest validators in time (yielding a polka of stake
\texttt{≥\ n\ −\ f\ ≥\ Q}), and similarly for precommits. The block
finalizes. Round-robin proposer selection (§4.2) guarantees an honest
proposer is selected in at most \texttt{f\ +\ 1} consecutive rounds,
well within \texttt{MAX\_ROUND\ =\ 100}. ∎

\hypertarget{48-message-complexity}{%
\subsubsection{4.8 Message Complexity}\label{48-message-complexity}}

Per round, each validator broadcasts one \texttt{Prevote} and one
\texttt{Precommit}. The proposer additionally broadcasts one
\texttt{Proposal}. With \texttt{n} validators, the per-round originating
messages are exactly \texttt{2n\ +\ 1} (one proposal, \texttt{n}
prevotes, \texttt{n} precommits). Under all-to-all delivery via
gossipsub mesh, total bytes transferred scale as \texttt{O(n²)} worst
case; with optimal gossip the total amortizes to \texttt{O(n\ log\ n)}.

For \texttt{n\ =\ 4} (current mainnet, \texttt{Q\ =\ 3}), per-round
messages are \texttt{2\ ×\ 4\ +\ 1\ =\ 9}. For the maximum active-set
size \texttt{n\ =\ 21} (\texttt{Q\ =\ 15}), per-round messages are
\texttt{2\ ×\ 21\ +\ 1\ =\ 43}. Both regimes are well within
libp2p\textquotesingle s gossip throughput envelope.

\hypertarget{49-bft-gate-relax}{%
\subsubsection{4.9 BFT Gate Relax}\label{49-bft-gate-relax}}

Pre-fork, the protocol requires the entire active set online
(\texttt{active\ =\ n}) to make consensus progress. Post
\texttt{BFT\_GATE\_RELAX\_HEIGHT}, the requirement relaxes to
\texttt{active\ ≥\ ⌈2n/3⌉}, providing a one-jail tolerance margin for
\texttt{n\ =\ 4} (allowing the chain to advance with three of four
signers when one is jailed). The gate is operator-controlled via the
\texttt{BFT\_GATE\_RELAX\_HEIGHT} environment variable; default is
\texttt{u64::MAX} (disabled).

\begin{center}\rule{0.5\linewidth}{0.5pt}\end{center}

\hypertarget{5-execution-layer}{%
\subsection{5. Execution Layer}\label{5-execution-layer}}

The execution layer implements \texttt{STF:\ 𝒮\ ×\ ℬ\ →\ 𝒮\ ∪\ \{⊥\}}.
It dispatches transactions to one of two rails --- \emph{native} or
\emph{EVM} --- based on the transaction\textquotesingle s
\texttt{to\_address}, applies side effects to the canonical state, and
computes the resulting state root.

\hypertarget{51-pipeline}{%
\subsubsection{5.1 Pipeline}\label{51-pipeline}}

\begin{Shaded}
\begin{verbatim}
Block B (txs in order)
  |
  v
Pass 1: signature verification (parallelizable per tx)
  |
  v
Pass 2: nonce + balance check
  |
  v
Fee deduct + ceil/floor burn split (burn debited but never credited)
  |
  v
Routing by tx.to_address:
  +- 0x...0000 --> Native TokenOp
  +- 0x...0100 --> Native StakingOp
  +- 0x...0002 --> System / Treasury
  +- other     --> EVM dispatch via revm 37
  |
  v
State mutation (sequential by B.txs order)
  |
  v
Pass 3: block subsidy (1 SRX -> PROTOCOL_TREASURY,
        accrued pro-rata to precommit signers of B_{h-1})
  |
  v
Trie commit: recompute root R from S'
  |
  v
If h >= STATE_ROOT_FORK_HEIGHT and R != B.header.state_root: REJECT BLOCK
Otherwise commit S' to chain.db
\end{verbatim}
\end{Shaded}

\textbf{Figure 3.} Execution pipeline. Each transaction passes through
signature verification, fee handling, rail dispatch, and state mutation.
The block subsidy is applied once per block at the end; the state trie
is recomputed, and the resulting root is checked against the
proposer\textquotesingle s claim. Pass 1 (signature verification) is
independently parallelizable per transaction; Pass 2 onwards is strictly
sequential by \texttt{B.txs} order (§5.2).

\hypertarget{52-determinism-axioms}{%
\subsubsection{5.2 Determinism Axioms}\label{52-determinism-axioms}}

For \texttt{STF} to satisfy the SMR Property (§2.4), every
implementation must be deterministic in three senses:

\textbf{(D1) Operation order.} Transactions are applied in
\texttt{B.txs} order. No reordering, no parallel-execution-with-merge
that produces a different commit order.

\textbf{(D2) State read consistency.} A transaction\textquotesingle s
state reads at its execution point reflect all prior transactions in the
same block. No snapshot isolation, no read-after-write divergence.

\textbf{(D3) Floating-point freedom.} No state mutation depends on
IEEE-754 arithmetic. All calculations are integer or fixed-point.
Sentrix uses \texttt{u64} sentri throughout; SRX-denominated
computations are integer multiples or quotients (e.g.
\texttt{total\_fee.div\_ceil(2)} for the burn share, §8.3).

These axioms hold for the current sequential implementation. They are
also the constraints any future parallel-execution scheme (§5.6) must
satisfy: a parallel execution is valid if and only if it produces a
state indistinguishable from the sequential execution defined by
\texttt{B.txs} order.

\hypertarget{53-state-transition-function}{%
\subsubsection{5.3 State Transition
Function}\label{53-state-transition-function}}

\begin{verbatim}
STF(S, B):
  // Pass 0: header validation
  1.  verify B.header.previous_hash == H(B_{h-1})
  2.  verify B.header.proposer == propose(B.header.index, r)
                                  for some round r ≤ MAX_ROUND
                                  (round determined by justification)
  3.  verify B.header.justification is a valid Q-precommit set on B_{h-1}
  4.  S' ← S

  // Pass 1: transaction signature verification (parallelizable)
  5.  for each tx in B.txs:
      a. verify ECDSA(tx.signature, canonical_signing_payload(tx),
                       tx.public_key)
      b. verify keccak256(tx.public_key)[12:] == tx.from_address

  // Pass 2: sequential apply (D1)
  6.  for each tx in B.txs (in order):
      a. if S'[tx.from].nonce ≠ tx.nonce: skip
      b. if tx.fee < MIN_TX_FEE: skip
      c. if S'[tx.from].balance < tx.fee + tx.amount: skip
      d. S'[tx.from].balance -= tx.fee
      e. // burn share is debited but never credited;
         // total_burned is tracked in chain accounting
         total_burned_srx += tx.fee.div_ceil(2)
      f. dispatch tx based on tx.to_address (Table 5.5):
         - TOKEN_OP_ADDRESS  (0x...0000) → native TokenOp
         - STAKING_ADDRESS   (0x...0100) → native StakingOp
         - PROTOCOL_TREASURY (0x...0002) → system/claim path
         - any other         → EVM dispatch via revm
      g. if dispatch fails:
           - revert state changes for this tx (step f side-effects)
           - keep fee deduction (step d) and burn debit (step e)
      h. S'[tx.from].nonce += 1

  // Pass 3: validator fee credit (per block, not per tx)
  7.  let total_fee = Σ tx.fee for tx in B.txs (admissible)
  8.  let validator_share = total_fee - total_fee.div_ceil(2)
                          // = floor(total_fee / 2)
  9.  S'[B.header.proposer].balance += validator_share

  // Pass 4: block subsidy
  10. let subsidy = block_reward(B.header.index)
                  = max(BLOCK_REWARD ÷ 2^era, 0) where
                    era = ⌊B.header.index / HALVING_INTERVAL_V2⌋
  11. S'[PROTOCOL_TREASURY].balance += subsidy
  12. for each precommitter v in B_{h-1}.justification.precommits:
        S'.staking.pending_rewards[v] += subsidy × s(v) / Σ s(precommitters)
                                         × (1 − commission(v) for delegators)

  // Pass 5: trie commit
  13. recompute root R from S' (binary sparse Merkle, depth 256)
  14. if h ≥ STATE_ROOT_FORK_HEIGHT and R ≠ B.header.state_root:
        return ⊥
  15. return S'
\end{verbatim}

\texttt{STF} is deterministic by D1--D3 (§5.2), so any honest validator
computes the same \texttt{S\textquotesingle{}} from the same
\texttt{(S,\ B)}. A divergence in step 14 indicates either (a) a
Byzantine proposer who proposed a block whose claimed state root does
not match correct execution, or (b) a bug in the local executor.

The state-root check is gated by
\texttt{STATE\_ROOT\_FORK\_HEIGHT\ =\ 100\_000}. Pre-fork blocks are not
committed to a state root in their header; the chain\textquotesingle s
bytewise determinism is the only consensus binding for those blocks.
Post-fork, the state root is part of the block hash and any divergence
aborts apply.

\hypertarget{54-native-rail}{%
\subsubsection{5.4 Native Rail}\label{54-native-rail}}

The native rail dispatches three operation classes based on
\texttt{to\_address}:

\textbf{TokenOp} (\texttt{TOKEN\_OP\_ADDRESS\ =\ 0x00…0000}) Encoded as
canonical JSON \texttt{\{"op":\ "...",\ ...\}} in \texttt{tx.data}.
Variants: \texttt{Deploy}, \texttt{Transfer}, \texttt{Burn},
\texttt{Mint}, \texttt{Approve} (SRC-20 fungibles). Variants
\texttt{DeployNft}, \texttt{MintNft}, \texttt{TransferNft},
\texttt{BurnNft}, \texttt{Deploy1155}, \texttt{Mint1155},
\texttt{Transfer1155}, \texttt{Burn1155} (SRC-721 / SRC-1155 NFTs) are
wire-format-stable but dispatch is gated by
\texttt{NFT\_TOKENOP\_HEIGHT\ =\ u64::MAX} (disabled) pending Pass-2
storage layer.

\textbf{StakingOp} (\texttt{STAKING\_ADDRESS\ =\ 0x00…0100}) Encoded as
canonical JSON. Variants: \texttt{RegisterValidator}, \texttt{Delegate},
\texttt{Undelegate}, \texttt{Redelegate}, \texttt{ClaimRewards},
\texttt{Unjail}, \texttt{AddSelfStake}, \texttt{SubmitEvidence},
\texttt{JailEvidenceBundle}. The last is gated by
\texttt{JAIL\_CONSENSUS\_HEIGHT\ =\ u64::MAX} (consensus-computed jail
dispatch deferred pending non-determinism root cause fix).

\textbf{System} (\texttt{PROTOCOL\_TREASURY\ =\ 0x00…0002}) Treasury
operations (reward escrow drain), implemented inside the executor rather
than as a user-facing op enum.

The native rail bypasses revm entirely. Native ops cost
\texttt{MIN\_TX\_FEE\ =\ 10\_000\ sentri} flat, regardless of operation.
Native operations interoperate with the EVM rail through the canonical
state: an EVM contract may read native SRC-20 balances through a
precompile; native staking state is similarly readable.

\hypertarget{55-evm-rail}{%
\subsubsection{5.5 EVM Rail}\label{55-evm-rail}}

For any other \texttt{to\_address}, the dispatcher invokes
\texttt{revm\ 37} via the \texttt{sentrix-evm} adapter with a state
provider backed by the chain trie. Standard EVM semantics apply: gas
accounting follows EIP-1559 {[}13{]} (per-block \texttt{baseFeePerGas}
plus sender-set tip), opcode behaviour matches Ethereum mainnet with the
exception of chain ID (\texttt{7119} mainnet, \texttt{7120} testnet).
Standard tooling (Foundry, Hardhat, MetaMask, ethers, viem)
interoperates without modification.

EVM-side \texttt{eth\_sendRawTransaction} enters a translation layer
that wraps the EIP-155 RLP-encoded transaction into the canonical
Sentrix transaction form (§5.7) before mempool admission. EIP-155
signature verification happens before translation; gossip and
finalization are identical from there on.

EVM value transfers (a non-zero \texttt{tx.value} accompanying a
contract call) are gated by
\texttt{EVM\_VALUE\_TRANSFER\_HEIGHT\ =\ u64::MAX} (disabled) pending
the eager-write divergence investigation that motivated the v2.1.50
fork-gate. Operators flip the gate via the
\texttt{EVM\_VALUE\_TRANSFER\_HEIGHT} environment variable when Pass-2
EVM apply is verified consistent across the cluster.

Routing is summarized in:

\begin{longtable}[]{@{}lll@{}}
\toprule\noalign{}
\texttt{to\_address} sentinel & Constant name & Rail / Class \\
\midrule\noalign{}
\endhead
\bottomrule\noalign{}
\endlastfoot
\texttt{0x0000000000000000000000000000000000000000} &
\texttt{TOKEN\_OP\_ADDRESS} & Native TokenOp \\
\texttt{0x0000000000000000000000000000000000000100} &
\texttt{STAKING\_ADDRESS} & Native StakingOp \\
\texttt{0x0000000000000000000000000000000000000002} &
\texttt{PROTOCOL\_TREASURY} & System / Treasury \\
anything else & --- & EVM via revm \\
\end{longtable}

\hypertarget{56-forward-looking-parallel-execution}{%
\subsubsection{5.6 Forward-Looking Parallel
Execution}\label{56-forward-looking-parallel-execution}}

Execution within a block is currently strictly sequential (D1, §5.2).
The throughput ceiling imposed by the protocol is

\begin{verbatim}
TPS_max = MAX_TX_PER_BLOCK / BLOCK_TIME_SECS = 5_000 / 1 = 5_000 tx/s
\end{verbatim}

The achievable rate on any specific deployment is bounded by
single-threaded apply throughput (signature verification, nonce/balance
checks, dispatch, trie writes). Beyond the regime where sequential apply
saturates one CPU core, parallel execution becomes necessary.

The forward-looking parallel-execution model adopts an \emph{optimistic
concurrency control} scheme inspired by Block-STM {[}16{]} and
Monad\textquotesingle s pipelined execution model {[}18{]}:

\begin{enumerate}
\def\labelenumi{\arabic{enumi}.}
\tightlist
\item
  \textbf{Static dependency hints.} Each transaction declares (or the
  dispatcher infers) a \emph{read-set} and \emph{write-set}. For native
  ops the sets are derived from operation arguments (e.g.
  \texttt{Delegate.validator} is the only stake-registry write). For EVM
  dispatch the sets are over-approximated from prior-block traces.
\item
  \textbf{Independent-batch identification.} Build a directed dependency
  graph \texttt{G\ =\ (V,\ E)} where \texttt{tx\_i\ →\ tx\_j} iff
  \texttt{tx\_j}\textquotesingle s read-set intersects
  \texttt{tx\_i}\textquotesingle s write-set and
  \texttt{j\ \textgreater{}\ i} in block order. Compute root nodes of
  \texttt{G}; these execute in parallel. After each batch commits,
  update \texttt{G} and identify the next layer.
\item
  \textbf{Optimistic re-execution.} Transactions whose dependencies are
  over-approximated execute speculatively; on a write-conflict, abort
  and re-execute the conflicting transaction sequentially.
\item
  \textbf{Sequential commit.} Apply each transaction\textquotesingle s
  writes to canonical state in \texttt{B.txs} order regardless of
  execution order. This preserves D1 --- the \emph{commit} schedule
  equals the block order.
\end{enumerate}

The required protocol changes are: (a) declared read/write sets in
transaction metadata, (b) a speculative-execution-friendly state
provider, (c) a deterministic conflict-resolution rule. Implementation
is sequenced after the chain reaches a stable single-implementation
baseline; current production runs the sequential apply path, and no
consensus-layer changes are required to enable parallel execution later.

\hypertarget{57-transaction-format}{%
\subsubsection{5.7 Transaction Format}\label{57-transaction-format}}

A transaction is the unit of state change. Its canonical wire format is
(\texttt{crates/sentrix-primitives/src/transaction.rs}):

\begin{Shaded}
\begin{Highlighting}[]
\KeywordTok{pub} \KeywordTok{struct}\NormalTok{ Transaction }\OperatorTok{\{}
    \KeywordTok{pub}\NormalTok{ txid}\OperatorTok{:}         \DataTypeTok{String}\OperatorTok{,}      \CommentTok{// hex{-}encoded SHA{-}256 of canonical signing payload}
    \KeywordTok{pub}\NormalTok{ from\_address}\OperatorTok{:} \DataTypeTok{String}\OperatorTok{,}      \CommentTok{// hex{-}encoded 20{-}byte address (ECDSA{-}recovered)}
    \KeywordTok{pub}\NormalTok{ to\_address}\OperatorTok{:}   \DataTypeTok{String}\OperatorTok{,}      \CommentTok{// hex{-}encoded 20{-}byte address (routing key, §5.5)}
    \KeywordTok{pub}\NormalTok{ amount}\OperatorTok{:}       \DataTypeTok{u64}\OperatorTok{,}         \CommentTok{// sentri (1 SRX = 10⁸ sentri)}
    \KeywordTok{pub}\NormalTok{ fee}\OperatorTok{:}          \DataTypeTok{u64}\OperatorTok{,}         \CommentTok{// sentri, ≥ MIN\_TX\_FEE = 10\_000}
    \KeywordTok{pub}\NormalTok{ nonce}\OperatorTok{:}        \DataTypeTok{u64}\OperatorTok{,}         \CommentTok{// sender account sequence}
    \KeywordTok{pub}\NormalTok{ data}\OperatorTok{:}         \DataTypeTok{String}\OperatorTok{,}      \CommentTok{// empty for plain transfers;}
                                   \CommentTok{// canonical JSON for native ops;}
                                   \CommentTok{// EVM call payload for EVM dispatch}
    \KeywordTok{pub}\NormalTok{ timestamp}\OperatorTok{:}    \DataTypeTok{u64}\OperatorTok{,}         \CommentTok{// unix seconds (anti{-}replay window)}
    \KeywordTok{pub}\NormalTok{ chain\_id}\OperatorTok{:}     \DataTypeTok{u64}\OperatorTok{,}         \CommentTok{// 7119 mainnet, 7120 testnet}
    \KeywordTok{pub}\NormalTok{ signature}\OperatorTok{:}    \DataTypeTok{String}\OperatorTok{,}      \CommentTok{// hex secp256k1 ECDSA, 65 bytes (r, s, v)}
    \KeywordTok{pub}\NormalTok{ public\_key}\OperatorTok{:}   \DataTypeTok{String}\OperatorTok{,}      \CommentTok{// hex secp256k1 uncompressed, 65 bytes}
\OperatorTok{\}}
\end{Highlighting}
\end{Shaded}

Maximum transaction size is \texttt{MAX\_TX\_SIZE\ =\ 128\ KB}. The
mempool admits at most \texttt{MAX\_MEMPOOL\_SIZE\ =\ 10\_000}
transactions globally, with \texttt{MAX\_MEMPOOL\_PER\_SENDER\ =\ 100}
per-sender; transactions older than
\texttt{MEMPOOL\_MAX\_AGE\_SECS\ =\ 3\_600} are evicted.

The signed payload is the canonical JSON serialization of the eight
content fields (\texttt{amount}, \texttt{chain\_id}, \texttt{data},
\texttt{fee}, \texttt{from}, \texttt{nonce}, \texttt{timestamp},
\texttt{to}) in lexicographic key order. The payload is hashed with
SHA-256; the resulting 32-byte digest is signed with secp256k1 ECDSA.

EVM-side \texttt{eth\_sendRawTransaction} enters the translation layer
described in §5.5; native-side wallets that sign the canonical Sentrix
payload directly (e.g. \texttt{solux.sentriscloud.com}) submit to the
\texttt{sentrix\_sendTransaction} endpoint.

\begin{center}\rule{0.5\linewidth}{0.5pt}\end{center}

\hypertarget{6-state-layer}{%
\subsection{6. State Layer}\label{6-state-layer}}

\hypertarget{61-trie-commitment}{%
\subsubsection{6.1 Trie Commitment}\label{61-trie-commitment}}

Canonical state is committed via a \emph{binary sparse Merkle tree} with
depth 256 {[}9{]} (\texttt{crates/sentrix-trie}). Each leaf is keyed by
a 256-bit hash of the account address and storage slot. The root hash
\texttt{R(S)} is included in the block header as
\texttt{B.header.state\_root}.

For account \texttt{a}, the proof of inclusion is \texttt{O(log₂\ N)}
where \texttt{N} is the number of populated leaves. Light clients verify
any fact about \texttt{S} against \texttt{R(S)} with a logarithmic-size
Merkle path.

The state root commitment is gated by
\texttt{STATE\_ROOT\_FORK\_HEIGHT\ =\ 100\_000} (§5.3). Below this
height the chain advanced by bytewise consensus only; above it the state
root is part of the block hash and any divergence aborts apply.

\hypertarget{62-state-persistence}{%
\subsubsection{6.2 State Persistence}\label{62-state-persistence}}

State is persisted in MDBX {[}10{]} (\texttt{crates/sentrix-storage}), a
memory-mapped key-value store. Read amplification is bounded; write
amplification is constant-time per leaf update. State size grows with
the number of unique accounts, contract storage slots, and historical
block records.

A node maintains exactly one \texttt{chain.db} file. Recovery from
corruption (e.g. \texttt{MDBX\_MAP\_FULL} from disk pressure, or
page-layout divergence after a forensic backup operation) is performed
by acquiring a canonical \texttt{chain.db} copy from a healthy peer via
byte-level transfer; this is operationally documented (§11.4) and
produces bit-identical state regardless of the receiving
node\textquotesingle s prior state.

\hypertarget{63-light-clients}{%
\subsubsection{6.3 Light Clients}\label{63-light-clients}}

A light client tracks block headers but no full state. To verify a fact
\texttt{f} (e.g. account balance at height \texttt{h}):

\begin{enumerate}
\def\labelenumi{\arabic{enumi}.}
\tightlist
\item
  Acquire \texttt{B\_h.header.state\_root} through the BFT justification
  chain or a trusted weak-subjectivity checkpoint.
\item
  Request a Merkle proof from any full node for the leaf encoding
  \texttt{f} against root \texttt{R\ =\ B\_h.header.state\_root}.
\item
  Verify the proof locally; reject if the recomputed root does not
  match.
\end{enumerate}

This makes constrained-environment usage (mobile, browser) feasible
without trusting any specific full node. A formal weak-subjectivity
protocol --- including checkpoint distribution, refresh cadence, and
validator-set-rotation tracking --- is open work (§14).

\begin{center}\rule{0.5\linewidth}{0.5pt}\end{center}

\hypertarget{7-network-layer}{%
\subsection{7. Network Layer}\label{7-network-layer}}

\hypertarget{71-topology}{%
\subsubsection{7.1 Topology}\label{71-topology}}

Validators and full nodes form a libp2p mesh
(\texttt{crates/sentrix-network}). Peer discovery uses Kademlia DHT
{[}12{]} with seed peer addresses for bootstrap. Each validator
publishes a \emph{validator advertisement} record containing its current
libp2p multiaddrs; advertisements are gossiped on a dedicated topic and
persisted in the DHT, enabling validators to maintain direct routes for
BFT messages without external coordination.

\hypertarget{72-gossip-topics}{%
\subsubsection{7.2 Gossip Topics}\label{72-gossip-topics}}

\begin{longtable}[]{@{}lll@{}}
\toprule\noalign{}
Topic & Payload & Audience \\
\midrule\noalign{}
\endhead
\bottomrule\noalign{}
\endlastfoot
\texttt{sentrix/blocks/1} & Finalized blocks & All nodes; receivers
validate and apply via \texttt{STF} \\
\texttt{sentrix/txs/1} & User transactions & All nodes; forwarded until
included in a proposed block \\
\texttt{sentrix/bft/1} & Proposals, prevotes, precommits & Validators
only; full nodes ignore \\
\texttt{sentrix/validator-adverts/1} & Validator advertisements &
Validators; used to maintain direct routes \\
\end{longtable}

\hypertarget{73-bootstrap-sequence}{%
\subsubsection{7.3 Bootstrap Sequence}\label{73-bootstrap-sequence}}

\begin{verbatim}
1. Connect to seed peers (configured at startup).
2. Run Kademlia DHT walk to populate the peer routing table.
3. Request the chain head from peers; identify the head with the
   longest stake-weighted justification chain.
4. Sync block-by-block from a peer until the local height matches
   the network head.
5. Subscribe to gossip topics; if a validator, begin participating
   in BFT for the next round.
\end{verbatim}

A node behind the chain tip applies blocks linearly. A node ahead of the
tip cannot exist except through Byzantine equivocation; equivocation
evidence is itself a slashable offence (§9.1).

\hypertarget{74-transport}{%
\subsubsection{7.4 Transport}\label{74-transport}}

Cross-validator transport runs over public IPv4 with TLS-protected
libp2p Noise; an earlier WireGuard mesh deployment was retired
2026-04-30 in favour of public-IP transport. Validator hosts open one
inbound TCP port for libp2p (default 30303). RPC and WebSocket endpoints
are exposed via separate HTTP processes fronted by an edge reverse
proxy.

\begin{center}\rule{0.5\linewidth}{0.5pt}\end{center}

\hypertarget{8-tokenomics}{%
\subsection{8. Tokenomics}\label{8-tokenomics}}

\hypertarget{81-supply}{%
\subsubsection{8.1 Supply}\label{81-supply}}

The maximum supply is
\texttt{MAX\_SUPPLY\_V2\ =\ 315,000,000\ SRX\ =\ 3.15\ ×\ 10¹⁶\ sentri},
enforced post \texttt{TOKENOMICS\_V2\_HEIGHT}. Pre-fork the cap was
\texttt{MAX\_SUPPLY\_V1\ =\ 210M\ SRX}; the fork lifted the cap and
altered halving cadence. There is no mechanism (governance, hard fork,
or otherwise) by which the post-fork cap changes.

\hypertarget{82-block-reward-and-halving}{%
\subsubsection{8.2 Block Reward and
Halving}\label{82-block-reward-and-halving}}

The base block reward is
\texttt{BLOCK\_REWARD\ =\ 1\ SRX\ =\ 100\_000\_000\ sentri}. The reward
halves every \texttt{HALVING\_INTERVAL\_V2\ =\ 126\_000\_000} blocks
(\textasciitilde4 years at 1 s blocks, modeled on
Bitcoin\textquotesingle s halving cadence). Pre-fork the interval was
\texttt{HALVING\_INTERVAL\_V1\ =\ 42\_000\_000} blocks
(\textasciitilde1.33 years). Formally:

\begin{verbatim}
era(h)         = ⌊h / HALVING_INTERVAL_V2⌋  for h ≥ TOKENOMICS_V2_HEIGHT
block_reward(h) = BLOCK_REWARD × 2^{−era(h)}
\end{verbatim}

The disinflationary supply curve converges asymptotically to the cap;
combined with the 63M premine (§8.4), maximum supply approaches 315M
over a multi-decade horizon.

\hypertarget{83-fee-mechanism}{%
\subsubsection{8.3 Fee Mechanism}\label{83-fee-mechanism}}

Every transaction pays a fee. The fee model differs by rail:

\begin{itemize}
\tightlist
\item
  \textbf{Native rail:}
  \texttt{fee\ =\ MIN\_TX\_FEE\ =\ 10\_000\ sentri\ =\ 10⁻⁴\ SRX} flat,
  regardless of operation.
\item
  \textbf{EVM rail:} \texttt{fee\ =\ gas\_used\ ×\ (base\_fee\ +\ tip)}
  per EIP-1559. \texttt{base\_fee} adjusts each block to maintain target
  gas usage; \texttt{tip} is sender-set.
\end{itemize}

Fee disposition (per block, computed over
\texttt{total\_fee\ =\ Σ\ tx.fee} across admissible transactions):

\begin{verbatim}
burn_share      = total_fee.div_ceil(2)        // ceiling division
validator_share = total_fee − burn_share       // floor(total_fee / 2)
\end{verbatim}

The burn share is \emph{debited from senders but never credited
anywhere}; it leaves the supply entirely. The chain\textquotesingle s
\texttt{total\_burned\_srx} accounting tracks the cumulative burn for
telemetry (\texttt{/chain/info}). The validator share is credited
directly to the proposer\textquotesingle s balance
(\texttt{coinbase\_validator}), immediately spendable. For odd-valued
fees the burn receives the extra sentri (ceiling) --- this preserves a
clean integer split across all fees with no rounding loss.

Block subsidy is \emph{not} paid as a fee; it follows a separate path
(§8.5) so issuance enters circulation only when claimed.

\hypertarget{84-premine}{%
\subsubsection{8.4 Premine}\label{84-premine}}

A premine of \texttt{63M\ SRX\ =\ 20\%\ of\ MAX\_SUPPLY\_V2} was
allocated at genesis across four accounts. Allocations are public and
verifiable on-chain:

\begin{longtable}[]{@{}lll@{}}
\toprule\noalign{}
Role & Amount & Address \\
\midrule\noalign{}
\endhead
\bottomrule\noalign{}
\endlastfoot
Founder & 21M SRX &
\texttt{0x5b5b06688dcdbe532353ac610aaff41af825279d} \\
Early Validator & 10.5M SRX &
\texttt{0x328d56b8174697ef6c9e40e19b7663797e16fa47} \\
Ecosystem Fund & 21M SRX &
\texttt{0xeb70fdefd00fdb768dec06c478f450c351499f14} \\
Reserve & 10.5M SRX &
\texttt{0x2578cad17e3e56c2970a5b5eab45952439f5ba97} \\
\end{longtable}

The remaining 80\% (252M SRX) issues via block reward over the halving
curve. Operational sub-allocation policies (grant programs, listings,
vesting commitments) are outside this specification; refer to the public
tokenomics document at \texttt{sentrixchain.com/docs/tokenomics}.

\hypertarget{85-reward-routing}{%
\subsubsection{8.5 Reward Routing}\label{85-reward-routing}}

Block subsidy is \emph{not} paid directly to the proposer. It is minted
into a system address \texttt{PROTOCOL\_TREASURY\ =\ 0x00…0002} at the
end of each block (gated by \texttt{VOYAGER\_REWARD\_V2\_HEIGHT}), then
accrued pro-rata to the \emph{precommitters} of the parent block:

\begin{verbatim}
for v in B_{h-1}.justification.precommits:
    pending_rewards[v] += subsidy × s(v) / Σ s(precommitters)
\end{verbatim}

A delegator inherits its validator\textquotesingle s accrual minus the
validator\textquotesingle s published commission rate. Validators and
delegators drain \texttt{pending\_rewards} into spendable balance via an
explicit \texttt{ClaimRewards} staking operation.

This separation has two consequences:

\begin{enumerate}
\def\labelenumi{\arabic{enumi}.}
\tightlist
\item
  \textbf{Supply invariance.} New SRX enters circulation only when
  claimed, providing a clean accounting boundary between issuance and
  distribution. Circulating supply at height \texttt{h} is
  \texttt{total\_minted(h)\ −\ total\_burned(h)\ −\ Σ\ pending\_rewards}.
\item
  \textbf{Proposer revenue isolation.} The proposer\textquotesingle s
  fee share is real-time spendable, decoupling propose-incentive from
  sign-incentive.
\end{enumerate}

\begin{center}\rule{0.5\linewidth}{0.5pt}\end{center}

\hypertarget{9-validator-economics--slashing}{%
\subsection{9. Validator Economics \&
Slashing}\label{9-validator-economics--slashing}}

\hypertarget{91-slashing-conditions}{%
\subsubsection{9.1 Slashing Conditions}\label{91-slashing-conditions}}

\begin{longtable}[]{@{}llll@{}}
\toprule\noalign{}
Trigger & Evidence & Slash & Jail \\
\midrule\noalign{}
\endhead
\bottomrule\noalign{}
\endlastfoot
Equivocation (double-sign) & Two signed votes (prevote or precommit) on
conflicting blocks at the same \texttt{(h,\ r)} &
\texttt{DOUBLE\_SIGN\_SLASH\_BP\ =\ 2\_000} (20\% of self-stake) &
Permanent (tombstone) \\
Downtime & Signed
\texttt{\textless{}\ MIN\_SIGNED\_PER\_WINDOW\ =\ 4\_320} blocks within
a trailing \texttt{LIVENESS\_WINDOW\ =\ 14\_400} blocks
(\textasciitilde4 hours at 1 s) & \texttt{DOWNTIME\_SLASH\_BP\ =\ 10}
(0.1\%) & \texttt{DOWNTIME\_JAIL\_BLOCKS\ =\ 600} blocks
(\textasciitilde10 min) \\
\end{longtable}

Slashed SRX is destroyed (debited from the validator\textquotesingle s
stake but not credited anywhere --- same mechanism as fee burn, §8.3).
This preserves the supply invariant and avoids the perverse incentive of
validators benefiting from peer slashing.

The downtime threshold is permissive intentionally ---
\texttt{4\_320\ /\ 14\_400\ =\ 30\%} signed minimum --- so legitimate
causes (kernel reboot, brief network disruption) do not jail. Repeat
offenders accumulate jailings; their unbonded stake decays over time.

The consensus-computed jail dispatch
(\texttt{StakingOp::JailEvidenceBundle}) is gated by
\texttt{JAIL\_CONSENSUS\_HEIGHT\ =\ u64::MAX} (disabled). Manual jailing
--- operator submits a \texttt{Jail} admin operation against a divergent
or misbehaving validator --- remains the operational path until the
consensus dispatch is verified deterministic across the cluster.

\hypertarget{92-validator-revenue}{%
\subsubsection{9.2 Validator Revenue}\label{92-validator-revenue}}

For a validator with stake fraction \texttt{f\ =\ s(v)\ /\ Σ\ s(Vₐ)} and
commission rate \texttt{c}:

\begin{verbatim}
expected_revenue_per_epoch ≈ EPOCH_LENGTH × subsidy × f
                             × (1 − fraction_delegated × (1 − c))
                          + 0.5 × f × epoch_fees
\end{verbatim}

The first term (signing revenue) accrues into \texttt{pending\_rewards}
and is escrowed until the validator and its delegators claim. The second
term (proposing revenue) credits to the proposer\textquotesingle s
spendable balance directly each block. Under round-robin proposer
selection (§4.2), each active validator proposes exactly \texttt{1/n} of
blocks in expectation per epoch.

\hypertarget{93-liveness-penalty}{%
\subsubsection{9.3 Liveness Penalty}\label{93-liveness-penalty}}

A validator that signs fewer than \texttt{MIN\_SIGNED\_PER\_WINDOW}
blocks within the trailing \texttt{LIVENESS\_WINDOW} is jailed: removed
from the active set and slashed \texttt{DOWNTIME\_SLASH\_BP\ /\ 10000}
of self-stake. Re-entry requires an explicit \texttt{Unjail} transaction
after \texttt{DOWNTIME\_JAIL\_BLOCKS}. If slashing dropped self-stake
below the genesis-configured minimum, \texttt{AddSelfStake} (§4.1) is
required before re-entry.

\begin{center}\rule{0.5\linewidth}{0.5pt}\end{center}

\hypertarget{10-performance-model}{%
\subsection{10. Performance Model}\label{10-performance-model}}

\hypertarget{101-throughput}{%
\subsubsection{10.1 Throughput}\label{101-throughput}}

Per-block transaction capacity is bounded by
\texttt{MAX\_TX\_PER\_BLOCK\ =\ 5\_000}. With block time
\texttt{BLOCK\_TIME\_SECS\ =\ 1}, the protocol-imposed throughput
ceiling is

\begin{verbatim}
TPS_max = MAX_TX_PER_BLOCK / BLOCK_TIME_SECS = 5_000 tx/s
\end{verbatim}

The achievable throughput in any specific deployment is bounded below
this by:

\begin{verbatim}
TPS_actual = min(TPS_max, TPS_exec, TPS_net)
\end{verbatim}

where:

\begin{itemize}
\tightlist
\item
  \texttt{TPS\_exec} is the rate at which a single validator can verify,
  dispatch, and apply transactions sequentially (D1, §5.2). Sequential
  apply is the binding constraint; the rate depends on the transaction
  mix (native ops are cheaper than EVM calls; EVM cost scales with
  opcode complexity), the validator\textquotesingle s CPU, and trie I/O
  latency. No public benchmark numbers are codified in this document;
  the prototype benchmark framework (§12) is the path to producing them.
\item
  \texttt{TPS\_net} is the rate at which gossipsub can propagate the
  proposal and votes within \texttt{BLOCK\_TIME\_SECS}. It is bounded by
  \texttt{bandwidth\ ×\ BLOCK\_TIME\ /\ \textbar{}B\textbar{}}. For
  \texttt{MAX\_TX\_PER\_BLOCK\ =\ 5\_000} at typical transaction sizes,
  this is well within standard deployment bandwidth.
\end{itemize}

For target operating regions (\texttt{n\ =\ 4} to \texttt{n\ =\ 21}),
\texttt{TPS\_exec} is the binding constraint. Future parallel execution
(§5.6) is the path to relaxing it.

\hypertarget{102-latency}{%
\subsubsection{10.2 Latency}\label{102-latency}}

Per-block latency in the happy path decomposes into:

\begin{verbatim}
T_block = T_propose + T_prevote + T_precommit + T_apply
\end{verbatim}

Each of the first three terms is \emph{bounded above} by its phase
timeout (§4.2) --- \texttt{propose\_timeout(0)\ =\ 20\_000\ ms},
\texttt{prevote\_timeout(0)\ =\ 12\_000\ ms},
\texttt{precommit\_timeout(0)\ =\ 12\_000\ ms} --- but in normal
operation completes far below the timeout. Empirically, on a healthy
mesh, the chain produces blocks at the 1 s target, implying typical
phase durations on the order of tens to low-hundreds of milliseconds.
\texttt{T\_apply} is the time to verify signatures, dispatch
transactions, mutate state, and commit the trie root.

End-to-end latency observed by a transaction submitter is

\begin{verbatim}
T_user = T_mempool + T_block_inclusion + T_finalization
       ≤ 1 round   + 1 block            + 1 block
\end{verbatim}

where \texttt{T\_mempool} is the time until a proposer drains the
transaction. Under non-adversarial conditions the proposer sees the
transaction before composing the next block, and \texttt{T\_user} is
dominated by \texttt{T\_block}.

Round-skip extends end-to-end latency proportionally. After \texttt{r}
failed rounds at a single height, the cumulative wall-clock cost is
bounded by
\texttt{Σ\ (propose\_timeout(i)\ +\ prevote\_timeout(i)\ +\ precommit\_timeout(i))}
for \texttt{i\ ∈\ {[}0,\ r{]}}, which saturates at
\texttt{3\ ×\ 30\_000\ ms\ =\ 90\ s} per round once timeouts hit the
cap.

\hypertarget{103-message-complexity}{%
\subsubsection{10.3 Message Complexity}\label{103-message-complexity}}

Per round, \texttt{2n\ +\ 1} consensus messages are originated (one
proposal, \texttt{n} prevotes, \texttt{n} precommits). Total bytes
transferred across the mesh scale as \texttt{O(n²)} worst case for naive
all-to-all delivery; with optimal gossipsub mesh \texttt{O(n\ log\ n)}.

\begin{longtable}[]{@{}lll@{}}
\toprule\noalign{}
\texttt{n} & Originating messages per round & Worst-case mesh
transfer \\
\midrule\noalign{}
\endhead
\bottomrule\noalign{}
\endlastfoot
4 & 9 & \texttt{O(16)} \\
21 & 43 & \texttt{O(441)} \\
100 & 201 & \texttt{O(10⁴)} \\
\end{longtable}

For target operating regions (\texttt{n\ ≤\ 21}), per-round message
volume is small and well within libp2p\textquotesingle s gossip
throughput envelope.

\hypertarget{104-scalability-bounds}{%
\subsubsection{10.4 Scalability Bounds}\label{104-scalability-bounds}}

The single-shard architecture imposes:

\begin{itemize}
\tightlist
\item
  \textbf{State growth} is linear in unique accounts and contract
  storage slots. The trie\textquotesingle s \texttt{O(log\ N)} proof
  size and \texttt{O(log\ N)} update cost mean state size growth is the
  limiting factor for long-horizon node disk requirements.
\item
  \textbf{Bandwidth per validator} is
  \texttt{O(n\ ×\ \textbar{}B\textbar{})} per block due to gossip
  fan-out.
\item
  \textbf{Compute per block} is \texttt{O(\textbar{}B\textbar{})}
  sequential apply; future parallel apply (§5.6) reduces this to
  \texttt{O(\textbar{}B\textbar{}\ /\ p)} for \texttt{p}-way
  parallelism.
\item
  \textbf{Validator-set size} is capped at
  \texttt{MAX\_ACTIVE\_VALIDATORS\ =\ 21} by the staking module. Beyond
  this, the active set rotates by stake ranking; candidates outside the
  top 21 do not participate in consensus.
\end{itemize}

These bounds frame the protocol\textquotesingle s operating envelope:
Sentrix is sized for retail-grade settlement at validator counts in the
low tens, not for thousand-validator general-purpose computation.
Sharding and L2 rollups are out of scope for the current specification.

\begin{center}\rule{0.5\linewidth}{0.5pt}\end{center}

\hypertarget{11-failure-handling}{%
\subsection{11. Failure Handling}\label{11-failure-handling}}

This section specifies behaviour under four failure classes that the
production deployment has encountered.

\hypertarget{111-network-partition}{%
\subsubsection{11.1 Network Partition}\label{111-network-partition}}

\textbf{Scenario.} The validator set splits into two groups by transient
network failure. Group \texttt{A} has stake fraction \texttt{α}, group
\texttt{B} has \texttt{1\ −\ α}.

\textbf{Protocol behaviour:}

\begin{itemize}
\tightlist
\item
  If \texttt{α\ \textless{}\ 2/3} and
  \texttt{1\ −\ α\ \textless{}\ 2/3}: neither group reaches \texttt{Q}.
  Both halt at the partition height. Round-skip (§4.4) extends timeouts
  up to the 30 s cap but cannot complete consensus. After
  \texttt{MAX\_ROUND\ =\ 100} consecutive failed rounds, the chain is
  considered stalled and operator intervention is required.
\item
  If \texttt{α\ ≥\ 2/3}: group \texttt{A} continues finalizing blocks;
  group \texttt{B} halts.
\item
  Recovery on heal: the minority group observes the longer
  stake-weighted justification chain from majority peers, validates it
  locally via the SMR property (§2.4), and rejoins by replaying
  canonical blocks.
\end{itemize}

\textbf{Operator action:} none required if partition heals naturally. If
the partition is permanent (e.g. a validator host is irrecoverable), the
operator coordinates a chain-state rsync from a canonical peer (§11.4).

\hypertarget{112-leader-equivocation}{%
\subsubsection{11.2 Leader Equivocation}\label{112-leader-equivocation}}

\textbf{Scenario.} A Byzantine proposer signs two different proposals at
the same \texttt{(h,\ r)} --- \texttt{Proposal(h,\ r,\ B)} and
\texttt{Proposal(h,\ r,\ B′)} with \texttt{H(B)\ ≠\ H(B′)}.

\textbf{Protocol behaviour:}

\begin{itemize}
\tightlist
\item
  The two proposals propagate to different subsets of validators
  (Byzantine adversary\textquotesingle s choice). Honest validators
  receiving only one proposal prevote for it; honest validators
  receiving both prevote \texttt{⊥}.
\item
  Safety holds (Theorem 1, §4.6): no two honest validators precommit
  conflicting blocks because the polka condition cannot simultaneously
  hold for both \texttt{H(B)} and \texttt{H(B′)} when ≥ \texttt{n\ −\ f}
  honest stake exists.
\item
  The two signed proposals constitute \emph{equivocation evidence}. Any
  honest node that observes both gossips a
  \texttt{StakingOp::SubmitEvidence} transaction containing both
  signatures.
\end{itemize}

\textbf{Slashing.} On inclusion of valid evidence, the proposer is
slashed \texttt{DOUBLE\_SIGN\_SLASH\_BP\ /\ 10\_000\ =\ 20\%} of
self-stake and tombstoned (permanent ban from re-entry).

\hypertarget{113-validator-downtime}{%
\subsubsection{11.3 Validator Downtime}\label{113-validator-downtime}}

\textbf{Scenario.} A validator is offline (kernel reboot, network
outage, hardware failure) for some duration.

\textbf{Protocol behaviour:}

\begin{itemize}
\tightlist
\item
  During downtime, the validator does not sign prevotes or precommits.
  Other validators continue if \texttt{≥\ Q} remain online.
\item
  The downtime is recorded against a moving window of
  \texttt{LIVENESS\_WINDOW\ =\ 14\_400} blocks. A validator that signs
  fewer than \texttt{MIN\_SIGNED\_PER\_WINDOW\ =\ 4\_320} blocks within
  the window is jailed at the next epoch boundary.
\item
  Jailing slashes \texttt{DOWNTIME\_SLASH\_BP\ /\ 10\_000\ =\ 0.1\%} of
  self-stake and removes the validator from the active set for
  \texttt{DOWNTIME\_JAIL\_BLOCKS\ =\ 600} blocks (\textasciitilde10
  minutes at 1 s blocks).
\item
  After the jail duration, the validator may submit
  \texttt{StakingOp::Unjail} to re-enter the active set, subject to
  having sufficient self-stake; if slashing dropped them below the
  genesis-configured minimum, \texttt{StakingOp::AddSelfStake} is
  required first.
\end{itemize}

\textbf{Cluster effect.} If \texttt{\textgreater{}\ f} validators are
simultaneously offline, the chain stalls:
\texttt{n\ −\ f\ −\ offline\ \textless{}\ Q}. The remaining online
validators continue proposing but cannot finalize. Recovery is automatic
on enough validators returning online. The
\texttt{BFT\_GATE\_RELAX\_HEIGHT} fork (§4.9) widens the jail-cascade
margin once activated.

\hypertarget{114-chain-recovery-after-partition}{%
\subsubsection{11.4 Chain Recovery After
Partition}\label{114-chain-recovery-after-partition}}

\textbf{Scenario.} A partition heals; a previously-minority partition
has stale state, or a node\textquotesingle s local \texttt{chain.db} has
diverged at the byte level (e.g. due to a hard-fork mis-application or a
misdirected forensic backup operation).

\textbf{Protocol behaviour:}

\begin{enumerate}
\def\labelenumi{\arabic{enumi}.}
\tightlist
\item
  The stale node\textquotesingle s BFT engine detects that peers report
  a higher finalized height.
\item
  The node enters \emph{block-sync mode}: requests blocks from peers in
  batches, validates each via \texttt{STF}, and applies them in order.
\item
  State convergence is guaranteed by the SMR property: applying the same
  blocks against the same prior state produces the same successor state.
\item
  Once the stale node reaches the network head, it rejoins consensus.
\end{enumerate}

\textbf{Operator-side recovery for byte-level chain.db divergence.} When
\texttt{STF} replay alone cannot bring a node to byte-parity (e.g.
because the stale node committed a divergent state under a fork
mis-apply, leaving its chain.db permanently off the canonical path), the
operator copies the canonical \texttt{chain.db} from a healthy peer:

\begin{verbatim}
operator# systemctl stop sentrix
operator# # PULL canonical → stale (NOT push)
operator# ssh canonical-peer 'tar -C /opt/sentrix/data -czf - chain.db' \
            | tar -C /opt/sentrix/data -xzf -
operator# chown -R sentriscloud:sentriscloud /opt/sentrix/data/chain.db
operator# systemctl start sentrix
\end{verbatim}

Direction of transfer matters: the canonical peer is the source; the
stale node is the destination. Post-recovery, MD5 parity must be
confirmed across the cluster
(\texttt{md5sum\ /opt/sentrix/data/chain.db/*.dat} on every validator
should produce identical hashes for the canonical files). Production
runbook details and incident-tested procedures live in operator
documentation.

\begin{center}\rule{0.5\linewidth}{0.5pt}\end{center}

\hypertarget{12-benchmark-framework}{%
\subsection{12. Benchmark Framework}\label{12-benchmark-framework}}

This section specifies a prototype benchmark for measuring
execution-engine throughput and latency. The implementation is intended
to live in the \texttt{tools/bench-tps/} crate of the workspace.

\hypertarget{121-sequential-engine-reference}{%
\subsubsection{12.1 Sequential Engine
(Reference)}\label{121-sequential-engine-reference}}

\begin{Shaded}
\begin{Highlighting}[]
\CommentTok{// Reference implementation matching the production STF Pass 2 exactly.}
\CommentTok{// Measures the baseline throughput against which parallel engines compare.}

\KeywordTok{fn}\NormalTok{ run\_sequential\_bench(state}\OperatorTok{:} \OperatorTok{\&}\KeywordTok{mut}\NormalTok{ State}\OperatorTok{,}\NormalTok{ txs}\OperatorTok{:} \OperatorTok{\&}\NormalTok{[Transaction]) }\OperatorTok{{-}\textgreater{}}\NormalTok{ Metrics }\OperatorTok{\{}
    \KeywordTok{let}\NormalTok{ t0 }\OperatorTok{=} \PreprocessorTok{Instant::}\NormalTok{now()}\OperatorTok{;}
    \KeywordTok{let} \KeywordTok{mut}\NormalTok{ applied }\OperatorTok{=} \DecValTok{0}\OperatorTok{;}
    \KeywordTok{let} \KeywordTok{mut}\NormalTok{ latencies }\OperatorTok{=} \DataTypeTok{Vec}\PreprocessorTok{::}\NormalTok{with\_capacity(txs}\OperatorTok{.}\NormalTok{len())}\OperatorTok{;}

    \ControlFlowTok{for}\NormalTok{ tx }\KeywordTok{in}\NormalTok{ txs }\OperatorTok{\{}
        \KeywordTok{let}\NormalTok{ t\_start }\OperatorTok{=} \PreprocessorTok{Instant::}\NormalTok{now()}\OperatorTok{;}

        \CommentTok{// Pass 1: signature verification}
        \ControlFlowTok{if} \OperatorTok{!}\NormalTok{verify\_signature(tx) }\OperatorTok{\{} \ControlFlowTok{continue}\OperatorTok{;} \OperatorTok{\}}

        \CommentTok{// Pass 2: nonce/balance check}
        \ControlFlowTok{if}\NormalTok{ state}\OperatorTok{.}\NormalTok{nonce(tx}\OperatorTok{.}\NormalTok{from) }\OperatorTok{!=}\NormalTok{ tx}\OperatorTok{.}\NormalTok{nonce }\OperatorTok{\{} \ControlFlowTok{continue}\OperatorTok{;} \OperatorTok{\}}
        \ControlFlowTok{if}\NormalTok{ state}\OperatorTok{.}\NormalTok{balance(tx}\OperatorTok{.}\NormalTok{from) }\OperatorTok{\textless{}}\NormalTok{ tx}\OperatorTok{.}\NormalTok{fee }\OperatorTok{+}\NormalTok{ tx}\OperatorTok{.}\NormalTok{amount }\OperatorTok{\{} \ControlFlowTok{continue}\OperatorTok{;} \OperatorTok{\}}

        \CommentTok{// Fee handling: ceil/floor split, burn debit (no credit)}
        \KeywordTok{let}\NormalTok{ burn\_share }\OperatorTok{=}\NormalTok{ tx}\OperatorTok{.}\NormalTok{fee}\OperatorTok{.}\NormalTok{div\_ceil(}\DecValTok{2}\NormalTok{)}\OperatorTok{;}
        \KeywordTok{let}\NormalTok{ validator\_share }\OperatorTok{=}\NormalTok{ tx}\OperatorTok{.}\NormalTok{fee }\OperatorTok{{-}}\NormalTok{ burn\_share}\OperatorTok{;}
\NormalTok{        state}\OperatorTok{.}\NormalTok{deduct(tx}\OperatorTok{.}\NormalTok{from}\OperatorTok{,}\NormalTok{ tx}\OperatorTok{.}\NormalTok{fee)}\OperatorTok{;}
\NormalTok{        state}\OperatorTok{.}\NormalTok{credit(PROPOSER\_ADDR}\OperatorTok{,}\NormalTok{ validator\_share)}\OperatorTok{;}
        \CommentTok{// burn\_share leaves supply entirely; track for accounting}
\NormalTok{        state}\OperatorTok{.}\NormalTok{total\_burned }\OperatorTok{+=}\NormalTok{ burn\_share}\OperatorTok{;}

        \CommentTok{// Dispatch}
        \ControlFlowTok{match}\NormalTok{ dispatch(tx) }\OperatorTok{\{}
            \ConstantTok{Ok}\NormalTok{(\_) }\OperatorTok{=\textgreater{}}\NormalTok{ apply(state}\OperatorTok{,}\NormalTok{ tx)}\OperatorTok{,}
            \ConstantTok{Err}\NormalTok{(\_) }\OperatorTok{=\textgreater{}} \OperatorTok{\{} \CommentTok{/* fee debit + burn stand; payload reverts */} \OperatorTok{\}}
        \OperatorTok{\}}
\NormalTok{        state}\OperatorTok{.}\NormalTok{increment\_nonce(tx}\OperatorTok{.}\NormalTok{from)}\OperatorTok{;}

\NormalTok{        latencies}\OperatorTok{.}\NormalTok{push(t\_start}\OperatorTok{.}\NormalTok{elapsed())}\OperatorTok{;}
\NormalTok{        applied }\OperatorTok{+=} \DecValTok{1}\OperatorTok{;}
    \OperatorTok{\}}

    \KeywordTok{let}\NormalTok{ elapsed }\OperatorTok{=}\NormalTok{ t0}\OperatorTok{.}\NormalTok{elapsed()}\OperatorTok{;}
\NormalTok{    Metrics }\OperatorTok{\{}
\NormalTok{        engine}\OperatorTok{:} \StringTok{"sequential"}\OperatorTok{,}
\NormalTok{        applied}\OperatorTok{,}
\NormalTok{        elapsed}\OperatorTok{,}
\NormalTok{        tps}\OperatorTok{:}\NormalTok{ applied }\KeywordTok{as} \DataTypeTok{f64} \OperatorTok{/}\NormalTok{ elapsed}\OperatorTok{.}\NormalTok{as\_secs\_f64()}\OperatorTok{,}
\NormalTok{        p50\_latency}\OperatorTok{:}\NormalTok{ percentile(}\OperatorTok{\&}\KeywordTok{mut}\NormalTok{ latencies}\OperatorTok{,} \DecValTok{50}\NormalTok{)}\OperatorTok{,}
\NormalTok{        p99\_latency}\OperatorTok{:}\NormalTok{ percentile(}\OperatorTok{\&}\KeywordTok{mut}\NormalTok{ latencies}\OperatorTok{,} \DecValTok{99}\NormalTok{)}\OperatorTok{,}
    \OperatorTok{\}}
\OperatorTok{\}}
\end{Highlighting}
\end{Shaded}

\hypertarget{122-batched-engine-optimistic-parallel}{%
\subsubsection{12.2 Batched Engine (Optimistic
Parallel)}\label{122-batched-engine-optimistic-parallel}}

\begin{Shaded}
\begin{Highlighting}[]
\CommentTok{// Batched engine: speculatively parallelizes transactions whose}
\CommentTok{// declared read/write sets do not conflict. On conflict, retries}
\CommentTok{// the offending transaction sequentially in commit order.}
\CommentTok{//}
\CommentTok{// Preserves D1 (commit order = block order) by accumulating writes}
\CommentTok{// in a per{-}tx WriteSet and merging into canonical state in order.}

\KeywordTok{fn}\NormalTok{ run\_batched\_bench(state}\OperatorTok{:} \OperatorTok{\&}\KeywordTok{mut}\NormalTok{ State}\OperatorTok{,}\NormalTok{ txs}\OperatorTok{:} \OperatorTok{\&}\NormalTok{[Transaction]}\OperatorTok{,}\NormalTok{ batch\_size}\OperatorTok{:} \DataTypeTok{usize}\NormalTok{) }\OperatorTok{{-}\textgreater{}}\NormalTok{ Metrics }\OperatorTok{\{}
    \KeywordTok{let}\NormalTok{ t0 }\OperatorTok{=} \PreprocessorTok{Instant::}\NormalTok{now()}\OperatorTok{;}
    \KeywordTok{let} \KeywordTok{mut}\NormalTok{ applied }\OperatorTok{=} \DecValTok{0}\OperatorTok{;}
    \KeywordTok{let} \KeywordTok{mut}\NormalTok{ latencies }\OperatorTok{=} \DataTypeTok{Vec}\PreprocessorTok{::}\NormalTok{with\_capacity(txs}\OperatorTok{.}\NormalTok{len())}\OperatorTok{;}

    \ControlFlowTok{for}\NormalTok{ chunk }\KeywordTok{in}\NormalTok{ txs}\OperatorTok{.}\NormalTok{chunks(batch\_size) }\OperatorTok{\{}
        \KeywordTok{let}\NormalTok{ graph }\OperatorTok{=}\NormalTok{ build\_dependency\_graph(chunk)}\OperatorTok{;}
        \KeywordTok{let}\NormalTok{ layers}\OperatorTok{:} \DataTypeTok{Vec}\OperatorTok{\textless{}}\DataTypeTok{Vec}\OperatorTok{\textless{}\&}\NormalTok{Transaction}\OperatorTok{\textgreater{}\textgreater{}} \OperatorTok{=}\NormalTok{ topological\_layers(}\OperatorTok{\&}\NormalTok{graph)}\OperatorTok{;}

        \ControlFlowTok{for}\NormalTok{ layer }\KeywordTok{in}\NormalTok{ layers }\OperatorTok{\{}
            \CommentTok{// Execute layer in parallel — each tx writes into its own WriteSet}
            \KeywordTok{let}\NormalTok{ writesets}\OperatorTok{:} \DataTypeTok{Vec}\OperatorTok{\textless{}}\NormalTok{WriteSet}\OperatorTok{\textgreater{}} \OperatorTok{=}\NormalTok{ layer}
                \OperatorTok{.}\NormalTok{par\_iter()}
                \OperatorTok{.}\NormalTok{map(}\OperatorTok{|}\NormalTok{tx}\OperatorTok{|} \OperatorTok{\{}
                    \KeywordTok{let}\NormalTok{ t\_start }\OperatorTok{=} \PreprocessorTok{Instant::}\NormalTok{now()}\OperatorTok{;}
                    \KeywordTok{let}\NormalTok{ ws }\OperatorTok{=}\NormalTok{ execute\_speculative(state}\OperatorTok{,}\NormalTok{ tx)}\OperatorTok{;}
\NormalTok{                    latencies}\OperatorTok{.}\NormalTok{push(t\_start}\OperatorTok{.}\NormalTok{elapsed())}\OperatorTok{;}
\NormalTok{                    ws}
                \OperatorTok{\}}\NormalTok{)}
                \OperatorTok{.}\NormalTok{collect()}\OperatorTok{;}

            \CommentTok{// Validate no write{-}write conflicts within layer}
            \ControlFlowTok{for}\NormalTok{ (i}\OperatorTok{,}\NormalTok{ ws\_i) }\KeywordTok{in}\NormalTok{ writesets}\OperatorTok{.}\NormalTok{iter()}\OperatorTok{.}\NormalTok{enumerate() }\OperatorTok{\{}
                \ControlFlowTok{for}\NormalTok{ ws\_j }\KeywordTok{in} \OperatorTok{\&}\NormalTok{writesets[}\OperatorTok{..}\NormalTok{i] }\OperatorTok{\{}
                    \ControlFlowTok{if}\NormalTok{ ws\_i}\OperatorTok{.}\NormalTok{conflicts(ws\_j) }\OperatorTok{\{}
                        \CommentTok{// Re{-}execute conflicting tx serially against committed}
                        \CommentTok{// state to maintain D1.}
                        \KeywordTok{let}\NormalTok{ tx\_i }\OperatorTok{=}\NormalTok{ layer[i]}\OperatorTok{;}
                        \KeywordTok{let}\NormalTok{ \_ws\_correct }\OperatorTok{=}\NormalTok{ execute\_speculative(state}\OperatorTok{,}\NormalTok{ tx\_i)}\OperatorTok{;}
                        \CommentTok{// Replace ws\_i with ws\_correct in the apply order...}
                    \OperatorTok{\}}
                \OperatorTok{\}}
            \OperatorTok{\}}

            \CommentTok{// Commit in block order}
            \ControlFlowTok{for}\NormalTok{ ws }\KeywordTok{in}\NormalTok{ writesets }\OperatorTok{\{}
\NormalTok{                state}\OperatorTok{.}\NormalTok{merge(ws)}\OperatorTok{;}
\NormalTok{                applied }\OperatorTok{+=} \DecValTok{1}\OperatorTok{;}
            \OperatorTok{\}}
        \OperatorTok{\}}
    \OperatorTok{\}}

    \KeywordTok{let}\NormalTok{ elapsed }\OperatorTok{=}\NormalTok{ t0}\OperatorTok{.}\NormalTok{elapsed()}\OperatorTok{;}
\NormalTok{    Metrics }\OperatorTok{\{}
\NormalTok{        engine}\OperatorTok{:} \StringTok{"batched"}\OperatorTok{,}
\NormalTok{        applied}\OperatorTok{,}
\NormalTok{        elapsed}\OperatorTok{,}
\NormalTok{        tps}\OperatorTok{:}\NormalTok{ applied }\KeywordTok{as} \DataTypeTok{f64} \OperatorTok{/}\NormalTok{ elapsed}\OperatorTok{.}\NormalTok{as\_secs\_f64()}\OperatorTok{,}
\NormalTok{        p50\_latency}\OperatorTok{:}\NormalTok{ percentile(}\OperatorTok{\&}\KeywordTok{mut}\NormalTok{ latencies}\OperatorTok{,} \DecValTok{50}\NormalTok{)}\OperatorTok{,}
\NormalTok{        p99\_latency}\OperatorTok{:}\NormalTok{ percentile(}\OperatorTok{\&}\KeywordTok{mut}\NormalTok{ latencies}\OperatorTok{,} \DecValTok{99}\NormalTok{)}\OperatorTok{,}
    \OperatorTok{\}}
\OperatorTok{\}}
\end{Highlighting}
\end{Shaded}

\hypertarget{123-metrics}{%
\subsubsection{12.3 Metrics}\label{123-metrics}}

Each run produces:

\begin{longtable}[]{@{}ll@{}}
\toprule\noalign{}
Metric & Definition \\
\midrule\noalign{}
\endhead
\bottomrule\noalign{}
\endlastfoot
\texttt{applied} & Number of transactions successfully applied (fee
debit + dispatch + state mutation) \\
\texttt{elapsed} & Wall-clock time from first tx received to last tx
committed \\
\texttt{tps} & \texttt{applied\ /\ elapsed} (transactions per second) \\
\texttt{p50\_latency} & Median per-tx execution time (verify + dispatch
+ apply) \\
\texttt{p99\_latency} & 99th percentile per-tx execution time \\
\texttt{conflict\_rate} & (Batched only) Fraction of
speculatively-executed txs aborted due to write-conflict \\
\end{longtable}

Comparative runs (sequential vs. batched at varying
\texttt{batch\_size}) characterize the speedup function and identify the
conflict regime where parallelism degrades into serial overhead.

\begin{center}\rule{0.5\linewidth}{0.5pt}\end{center}

\hypertarget{13-comparative-analysis}{%
\subsection{13. Comparative Analysis}\label{13-comparative-analysis}}

We compare Sentrix to four contemporary chains along three axes:
execution model, consensus design, and scalability approach.

\hypertarget{131-execution-model}{%
\subsubsection{13.1 Execution Model}\label{131-execution-model}}

\begin{longtable}[]{@{}llll@{}}
\toprule\noalign{}
Chain & Model & Determinism mechanism & Parallelism \\
\midrule\noalign{}
\endhead
\bottomrule\noalign{}
\endlastfoot
\textbf{Ethereum} & Sequential EVM (post-Merge) & Block-order serial
apply & None at L1 \\
\textbf{Solana} & Parallel SVM (Sealevel) & Account access lists
declared per tx & Native parallel; bounded by access-set declaration
accuracy \\
\textbf{Monad} & Optimistic parallel EVM & Block-STM-style speculative
execution + serial commit & Native parallel; conflict-resolved at
commit \\
\textbf{Polygon (PoS)} & Sequential EVM & Block-order serial apply &
None at the PoS chain layer \\
\textbf{Sentrix} & Sequential EVM + Native rail & Block-order serial
apply (D1, §5.2) & None today; optimistic parallel scheduled (§5.6) \\
\end{longtable}

The dominant axis distinguishing Sentrix is the \emph{native rail} ---
common operations (token issuance, staking, validator coordination)
bypass the EVM entirely and apply directly against canonical state. This
is similar in spirit to Cosmos SDK modules but co-resident with EVM
dispatch in the same node.

\hypertarget{132-consensus-design}{%
\subsubsection{13.2 Consensus Design}\label{132-consensus-design}}

\begin{longtable}[]{@{}lllll@{}}
\toprule\noalign{}
Chain & Consensus & Finality & Validator set & Per-round message
complexity \\
\midrule\noalign{}
\endhead
\bottomrule\noalign{}
\endlastfoot
\textbf{Ethereum} & Casper FFG + LMD-GHOST & Probabilistic (2 epochs ≈
12.8 min) & \textasciitilde1M (32 ETH per validator) & \texttt{O(n)} per
slot via committees \\
\textbf{Solana} & TowerBFT + Proof of History & Probabilistic ≈ 13s &
low thousands & \texttt{O(n)} per slot \\
\textbf{Monad} & MonadBFT (HotStuff-derivative) & Single-block &
Permissioned bootstrap → permissionless & \texttt{O(n)} per round
(threshold-aggregated) \\
\textbf{Polygon (PoS)} & Heimdall + Bor (Tendermint + Geth fork) &
\textasciitilde2 s nominal, \textasciitilde4 min Ethereum-checkpointed &
\textasciitilde100 & \texttt{O(n²)} per round \\
\textbf{Sentrix} & Voyager BFT (Tendermint-derivative) + DPoS &
Single-block & DPoS open, ranked by stake, capped at
\texttt{MAX\_ACTIVE\_VALIDATORS\ =\ 21} & \texttt{O(n²)} per round,
\texttt{O(n\ log\ n)} with optimal gossip \\
\end{longtable}

Sentrix is closest in lineage to Polygon\textquotesingle s PoS chain
(both are Tendermint-derivative BFT chains running an EVM rail). The
distinguishing choice in Sentrix is the native rail co-located with EVM
and the explicit coupling of DPoS validator selection to a single
round-robin proposer schedule. MonadBFT\textquotesingle s
threshold-aggregated O(n) view-change is a notable departure from the
Tendermint family Sentrix follows.

\hypertarget{133-scalability-approach}{%
\subsubsection{13.3 Scalability
Approach}\label{133-scalability-approach}}

\begin{longtable}[]{@{}llll@{}}
\toprule\noalign{}
Chain & Strategy & Bottleneck & Validator-set practical limit \\
\midrule\noalign{}
\endhead
\bottomrule\noalign{}
\endlastfoot
\textbf{Ethereum} & L2 rollups (Optimistic + ZK) & L1 data availability;
rollup batch size & L1 sustainable at \textasciitilde1M; L2 chains
compete for blob space \\
\textbf{Solana} & Vertical scaling (faster hardware, larger blocks) &
Network bandwidth + state I/O & low thousands \\
\textbf{Monad} & Pipelined parallel execution on a single shard &
Single-shard execution still bounded by hardware & HotStuff-class
\textasciitilde100 \\
\textbf{Polygon (PoS)} & Sidechain + zkEVM & Bridge security +
checkpoint cadence & \textasciitilde100 \\
\textbf{Sentrix} & Vertical (planned parallel exec, §5.6) & Sequential
apply + BFT message complexity & \texttt{MAX\_ACTIVE\_VALIDATORS\ =\ 21}
enforced \\
\end{longtable}

Sentrix\textquotesingle s scalability approach matches the
\emph{Tendermint-class} envelope: low-tens of validators, single-block
finality, vertical scaling within that bound. Sharding and L2 rollups
are out of scope for the current specification; the path forward through
§5.6\textquotesingle s parallel-execution model addresses the
sequential-apply bottleneck without changing consensus.

\hypertarget{134-posture-summary}{%
\subsubsection{13.4 Posture Summary}\label{134-posture-summary}}

Sentrix\textquotesingle s design is most accurately described as: a
single-shard EVM-compatible Tendermint BFT chain with a co-located
native operation rail and a deflationary capped-supply monetary policy.
None of these properties is novel in isolation; the contribution is the
combination, the determination to ship as a single small Rust binary,
and the operational simplicity of one process per validator host.

\begin{center}\rule{0.5\linewidth}{0.5pt}\end{center}

\hypertarget{14-open-problems}{%
\subsection{14. Open Problems}\label{14-open-problems}}

This specification leaves the following questions deliberately open.
They are listed honestly so independent reviewers can locate the
engineering frontier of the protocol.

\begin{enumerate}
\def\labelenumi{\arabic{enumi}.}
\tightlist
\item
  \textbf{Parallel execution determinism proof.} §5.6 describes the
  optimistic-concurrency model in outline. A rigorous proof that the
  proposed scheduler produces a state indistinguishable from sequential
  \texttt{B.txs}-order execution (D1, §5.2) is required before
  deployment.
\item
  \textbf{Light-client weak-subjectivity protocol.} §6.3 assumes a
  trusted checkpoint. A formal scheme --- including checkpoint
  distribution, refresh cadence ≤ unbonding period, and
  validator-set-rotation tracking --- is needed.
\item
  \textbf{Cross-rail atomicity.} A transaction spanning the native rail
  and the EVM rail (e.g. a contract call that triggers a native staking
  operation) currently uses a system gateway. A formal atomicity
  guarantee --- either both effects commit or neither --- is desirable.
\item
  \textbf{NFT TokenOp activation.} SRC-721 + SRC-1155 wire formats are
  stable in \texttt{crates/sentrix-primitives/src/transaction.rs}.
  Pass-2 dispatch + storage layer is not yet implemented; activation via
  \texttt{NFT\_TOKENOP\_HEIGHT} remains gated at \texttt{u64::MAX} until
  both ship.
\item
  \textbf{Consensus-computed jail.} \texttt{JAIL\_CONSENSUS\_HEIGHT} is
  gated at \texttt{u64::MAX} pending a non-determinism root cause fix;
  the bug fired twice on mainnet (2026-04-29, 2026-04-30) before the
  gate was returned to disabled. Manual jailing remains the operational
  path.
\item
  \textbf{EVM value-transfer fork-gate retirement.}
  \texttt{EVM\_VALUE\_TRANSFER\_HEIGHT} is gated at \texttt{u64::MAX}
  after a regression in v2.1.49 produced eager-write divergence. The
  fork-gate (introduced in v2.1.50) makes the new behaviour activatable
  on demand; a permanent retirement of the gate awaits a clean
  reproducer + fix.
\item
  \textbf{Multi-implementation diversity.} The protocol currently has a
  single Rust implementation. This specification at the level of detail
  given is the first step toward independent re-implementation; concrete
  client diversity is a long-horizon goal.
\item
  \textbf{Founder vesting on-chain.} Per §8.4, the Founder allocation
  has no on-chain vesting schedule; vesting is a public social
  commitment. A non-revocable linear schedule contract deployed via
  \texttt{SentrixSafe} is in the operational backlog.
\end{enumerate}

These open problems do not affect the safety or liveness of the current
production deployment. They represent the protocol\textquotesingle s
engineering frontier.

\begin{center}\rule{0.5\linewidth}{0.5pt}\end{center}

\hypertarget{15-references}{%
\subsection{15. References}\label{15-references}}

{[}1{]} Nakamoto, S. (2008). \emph{Bitcoin: A Peer-to-Peer Electronic
Cash System.}

{[}2{]} Buterin, V. (2014). \emph{Ethereum: A Next-Generation Smart
Contract and Decentralized Application Platform.}

{[}3{]} Buchman, E., Kwon, J. (2016). \emph{Tendermint: Byzantine Fault
Tolerance in the Age of Blockchains.}

{[}4{]} Bluealloy (2022). \emph{revm: Rust Ethereum Virtual Machine.}

{[}5{]} Yakovenko, A. (2017). \emph{Solana: A new architecture for a
high performance blockchain.}

{[}6{]} Kwon, J., Buchman, E. (2019). \emph{Cosmos: A Network of
Distributed Ledgers.}

{[}7{]} Fischer, M., Lynch, N., Paterson, M. (1985). \emph{Impossibility
of Distributed Consensus with One Faulty Process.} JACM 32(2).

{[}8{]} Castro, M., Liskov, B. (1999). \emph{Practical Byzantine Fault
Tolerance.} OSDI.

{[}9{]} Merkle, R. (1980). \emph{Protocols for Public Key
Cryptosystems.} IEEE S\&P.

{[}10{]} MDBX. \emph{Memory-mapped key-value store.}
\texttt{https://github.com/erthink/libmdbx}

{[}11{]} libp2p. \emph{Modular peer-to-peer networking stack.}
\texttt{https://libp2p.io}

{[}12{]} Maymounkov, P., Mazières, D. (2002). \emph{Kademlia: A
Peer-to-peer Information System Based on the XOR Metric.}

{[}13{]} Buterin, V. et al. (2019). \emph{EIP-1559: Fee market change
for ETH 1.0 chain.}

{[}14{]} Buterin, V., Griffith, V. (2017). \emph{Casper the Friendly
Finality Gadget.}

{[}15{]} Dwork, C., Lynch, N., Stockmeyer, L. (1988). \emph{Consensus in
the Presence of Partial Synchrony.} JACM 35(2).

{[}16{]} Gelashvili, R. et al. (2023). \emph{Block-STM: Scaling
Blockchain Execution by Turning Ordering Curse to a Performance
Blessing.} PPoPP.

{[}17{]} Yin, M. et al. (2019). \emph{HotStuff: BFT Consensus in the
Lens of Blockchain.} PODC.

{[}18{]} Monad Labs (2024). \emph{Monad: Parallelizing the EVM.}
Technical report.

\begin{center}\rule{0.5\linewidth}{0.5pt}\end{center}

\hypertarget{appendix-a--protocol-parameters}{%
\subsection{Appendix A --- Protocol
Parameters}\label{appendix-a--protocol-parameters}}

\begin{longtable}[]{@{}lll@{}}
\toprule\noalign{}
Parameter & Value & Source \\
\midrule\noalign{}
\endhead
\bottomrule\noalign{}
\endlastfoot
\texttt{BLOCK\_TIME\_SECS} & 1 &
\texttt{crates/sentrix-core/src/blockchain.rs} \\
\texttt{MAX\_TX\_PER\_BLOCK} & 5,000 &
\texttt{crates/sentrix-core/src/blockchain.rs} \\
\texttt{MAX\_TX\_SIZE} & 128 KB &
\texttt{crates/sentrix-core/src/mempool.rs} \\
\texttt{MAX\_MEMPOOL\_SIZE} & 10,000 &
\texttt{crates/sentrix-core/src/blockchain.rs} \\
\texttt{MAX\_MEMPOOL\_PER\_SENDER} & 100 &
\texttt{crates/sentrix-core/src/blockchain.rs} \\
\texttt{MEMPOOL\_MAX\_AGE\_SECS} & 3,600 &
\texttt{crates/sentrix-core/src/blockchain.rs} \\
\texttt{BLOCK\_REWARD} & 100,000,000 sentri (= 1 SRX) &
\texttt{crates/sentrix-core/src/blockchain.rs} \\
\texttt{MAX\_SUPPLY\_V2} & 315,000,000 SRX &
\texttt{crates/sentrix-core/src/blockchain.rs} \\
\texttt{HALVING\_INTERVAL\_V2} & 126,000,000 blocks &
\texttt{crates/sentrix-core/src/blockchain.rs} \\
\texttt{MIN\_TX\_FEE} & 10,000 sentri &
\texttt{crates/sentrix-primitives/src/transaction.rs} \\
\texttt{STATE\_ROOT\_FORK\_HEIGHT} & 100,000 &
\texttt{crates/sentrix-primitives/src/block.rs} \\
\texttt{MAX\_ACTIVE\_VALIDATORS} & 21 &
\texttt{crates/sentrix-staking/src/staking.rs} \\
\texttt{UNBONDING\_PERIOD} & 201,600 blocks &
\texttt{crates/sentrix-staking/src/staking.rs} \\
\texttt{EPOCH\_LENGTH} & 28,800 blocks (\textasciitilde8 hours) &
\texttt{crates/sentrix-staking/src/epoch.rs} \\
\texttt{LIVENESS\_WINDOW} & 14,400 blocks (\textasciitilde4 hours) &
\texttt{crates/sentrix-staking/src/slashing/liveness.rs} \\
\texttt{MIN\_SIGNED\_PER\_WINDOW} & 4,320 blocks &
\texttt{crates/sentrix-staking/src/slashing/liveness.rs} \\
\texttt{DOWNTIME\_JAIL\_BLOCKS} & 600 &
\texttt{crates/sentrix-staking/src/slashing/liveness.rs} \\
\texttt{DOWNTIME\_SLASH\_BP} & 10 (0.1\%) &
\texttt{crates/sentrix-staking/src/slashing/liveness.rs} \\
\texttt{DOUBLE\_SIGN\_SLASH\_BP} & 2,000 (20\%) &
\texttt{crates/sentrix-staking/src/slashing/double\_sign.rs} \\
\texttt{PROPOSE\_TIMEOUT\_MS} & 20,000 (cap 30,000) &
\texttt{crates/sentrix-bft/src/engine/timeouts.rs} \\
\texttt{PREVOTE\_TIMEOUT\_MS} & 12,000 (cap 30,000) &
\texttt{crates/sentrix-bft/src/engine/timeouts.rs} \\
\texttt{PRECOMMIT\_TIMEOUT\_MS} & 12,000 (cap 30,000) &
\texttt{crates/sentrix-bft/src/engine/timeouts.rs} \\
\texttt{TIMEOUT\_INCREMENT\_MS} & 1,000 (propose), 2,000 (votes) &
\texttt{crates/sentrix-bft/src/engine/timeouts.rs} \\
\texttt{MAX\_ROUND} & 100 &
\texttt{crates/sentrix-bft/src/engine/timeouts.rs} \\
\texttt{MAX\_TIMEOUT\_MS} & 30,000 &
\texttt{crates/sentrix-bft/src/engine/timeouts.rs} \\
\end{longtable}

\begin{longtable}[]{@{}lll@{}}
\toprule\noalign{}
Sentinel address & Value & Use \\
\midrule\noalign{}
\endhead
\bottomrule\noalign{}
\endlastfoot
\texttt{TOKEN\_OP\_ADDRESS} & \texttt{0x0000…0000} & Native TokenOp
routing \\
\texttt{STAKING\_ADDRESS} & \texttt{0x0000…0100} & Native StakingOp
routing \\
\texttt{PROTOCOL\_TREASURY} & \texttt{0x0000…0002} & System operations +
reward escrow \\
\end{longtable}

\begin{longtable}[]{@{}lll@{}}
\toprule\noalign{}
Fork-gate env var & Default & Effect \\
\midrule\noalign{}
\endhead
\bottomrule\noalign{}
\endlastfoot
\texttt{VOYAGER\_FORK\_HEIGHT} & \texttt{u64::MAX} & DPoS proposer
rotation + 3-phase BFT finality \\
\texttt{VOYAGER\_EVM\_HEIGHT} & \texttt{u64::MAX} & Embedded revm
runtime for \texttt{eth\_sendRawTransaction} \\
\texttt{VOYAGER\_REWARD\_V2\_HEIGHT} & \texttt{u64::MAX} & Coinbase
routes to PROTOCOL\_TREASURY; rewards claimed \\
\texttt{TOKENOMICS\_V2\_HEIGHT} & \texttt{u64::MAX} & 315M cap + 126M
halving (BTC-parity 4y at 1 s blocks) \\
\texttt{BFT\_GATE\_RELAX\_HEIGHT} & \texttt{u64::MAX} & BFT runs with
\texttt{active\ ≥\ ⌈2/3\ ×\ n⌉} instead of full mesh \\
\texttt{ADD\_SELF\_STAKE\_HEIGHT} & \texttt{u64::MAX} &
\texttt{StakingOp::AddSelfStake} dispatch active \\
\texttt{JAIL\_CONSENSUS\_HEIGHT} & \texttt{u64::MAX} &
Consensus-computed jail dispatch (currently disabled) \\
\texttt{NFT\_TOKENOP\_HEIGHT} & \texttt{u64::MAX} & SRC-721 + SRC-1155
dispatch (currently disabled) \\
\texttt{EVM\_VALUE\_TRANSFER\_HEIGHT} & \texttt{u64::MAX} & EVM tx.value
transfer plumbing (currently disabled) \\
\end{longtable}

Mainnet activation heights for these gates are operator-managed via
environment variables; default values (\texttt{u64::MAX}) mean disabled
until explicitly set.

\begin{center}\rule{0.5\linewidth}{0.5pt}\end{center}

\hypertarget{appendix-b--risk-disclosures}{%
\subsection{Appendix B --- Risk
Disclosures}\label{appendix-b--risk-disclosures}}

This appendix documents known risks. It is not exhaustive.

\textbf{Technical risks.}

\begin{itemize}
\tightlist
\item
  \emph{Single-implementation risk.} The chain has one Rust
  implementation. A bug in that implementation can affect the entire
  network until patched. Multiple-implementation diversity is not
  present.
\item
  \emph{Validator concentration during bootstrap.} The active validator
  set is small at this stage. Sustained Byzantine behavior by a
  supermajority is theoretically possible until the active set grows
  toward \texttt{MAX\_ACTIVE\_VALIDATORS\ =\ 21} with credibly
  independent operators.
\item
  \emph{Consensus-jail dispatch deferral.}
  \texttt{JAIL\_CONSENSUS\_HEIGHT} remains gated at \texttt{u64::MAX}
  pending a non-determinism root cause fix. Manual jailing remains
  operational.
\item
  \emph{EVM value-transfer fork-gate.}
  \texttt{EVM\_VALUE\_TRANSFER\_HEIGHT} remains gated at
  \texttt{u64::MAX} pending a clean reproducer + fix for the eager-write
  divergence regression that motivated the v2.1.50 gate.
\end{itemize}

\textbf{Economic risks.}

\begin{itemize}
\tightlist
\item
  \emph{Founder allocation is unilaterally controllable.} Until
  SentrixSafe expands to N-of-M and/or the Founder allocation is
  migrated to an on-chain vesting contract, the Founder address can move
  21M SRX at any time.
\item
  \emph{Bridge dependency for stablecoin counterparties.} Until a bridge
  protocol is deployed, Sentrix is islanded from the broader stablecoin
  economy.
\item
  \emph{No price discovery yet.} SRX is not currently trading on any
  exchange or DEX. The 315M cap is a protocol constant, not a market
  valuation.
\end{itemize}

\textbf{Operational risks.}

\begin{itemize}
\tightlist
\item
  \emph{Single-author bus factor.} Sentrix is solo-built. The
  author\textquotesingle s continued participation is not guaranteed by
  any contract. Long-term resilience depends on the codebase being open
  enough that other operators can run forks (BUSL → Apache 2.0
  transition after Change Date).
\item
  \emph{Infrastructure single points of failure.} The validator hosts,
  the explorer host, the DNS configuration, the documentation site ---
  all are operationally maintained by the author at this stage.
  Decentralization of these layers is a forward objective.
\end{itemize}

\textbf{Regulatory risks.}

\begin{itemize}
\tightlist
\item
  \emph{Securities-law classification.} SRX is intended as a utility
  token (gas, staking, governance). Whether any specific jurisdiction
  classifies it as a security depends on local law and how it is
  offered.
\item
  \emph{Cross-jurisdictional uncertainty.} Holders and validators in
  different jurisdictions face different regulatory regimes that may
  classify SRX or staking activity differently. Consult local counsel.
\end{itemize}

This list is honest, not exhaustive. The chain\textquotesingle s
behavior, contract, and history are public; readers should verify what
they care about against the on-chain record and the source code.

\begin{center}\rule{0.5\linewidth}{0.5pt}\end{center}

\hypertarget{appendix-c--legal-notice}{%
\subsection{Appendix C --- Legal
Notice}\label{appendix-c--legal-notice}}

This whitepaper is a description of a software protocol. It is not an
offering document, prospectus, investment solicitation, or financial
advice. SRX is a utility token used to pay transaction fees, secure the
chain through staking, and (in the future) participate in governance.
Acquiring or holding SRX entails risks --- technical, economic,
regulatory, and operational --- described in Appendix B.

The Sentrix Chain protocol is open-source under the Business Source
License 1.1, transitioning to Apache 2.0 after the Change Date specified
in the LICENSE file. Anyone may run a node, anyone may submit
transactions, anyone may operate as a validator subject to the
protocol-level requirements described in §4.1 and §9.

The author makes no representations about future SRX market price,
ecosystem adoption, partnership outcomes, or regulatory acceptance.

Readers are responsible for compliance with their local laws and
regulations regarding cryptocurrency holding, transaction, and validator
operation.

\begin{center}\rule{0.5\linewidth}{0.5pt}\end{center}

\hypertarget{about-the-author}{%
\subsection{About the Author}\label{about-the-author}}

\textbf{Satya Kwok} built Sentrix Chain solo in Rust. The
author\textquotesingle s prior work and engagement is publicly
observable on GitHub (\texttt{@satyakwok}) and through the
project\textquotesingle s open commit history at
\texttt{github.com/sentrix-labs/sentrix}. The author is contactable
through the channels listed at \texttt{sentrixchain.com}.

The decision to build Sentrix solo was deliberate: small teams ship
faster, decisions are clearer, and accountability is unambiguous. The
trade-off is the author\textquotesingle s bus factor (Appendix B). The
author is committed to operating Sentrix as durable infrastructure for
the indefinite future, but commits to no specific timeline beyond the
protocol-level guarantees codified in the chain itself.

\begin{center}\rule{0.5\linewidth}{0.5pt}\end{center}

\hypertarget{acknowledgments}{%
\subsection{Acknowledgments}\label{acknowledgments}}

Sentrix is the product of a single author\textquotesingle s work. The
author acknowledges the engineers who built the foundations on which
Sentrix is constructed --- the Ethereum core developers, the Cosmos and
Tendermint teams, the Bitcoin community, the Rust ecosystem, and the
open-source maintainers of the cryptographic and networking libraries
that make a project of this scope feasible for a single individual to
undertake.

The decision to build Sentrix in Rust, on the EVM, with Tendermint-style
BFT consensus, did not require inventing any of these components. It
required only their composition. This is how durable infrastructure is
built: not from new ideas, but from the careful arrangement of existing
ones.


\end{document}
